\begin{theorem}\label{night:P3K.119}
\StatementRecord{night:P3K.119}{proved}{}{\cref{night:P3K.116}; \cref{night:P3K.78}; \cref{night:P3F.35}; \cref{night:P3K.10} (the retained source, cap and integral local-system maps)}
Retain the original OS45 marked source and point line
\[
\begin{gathered}
R=t^2+a,\qquad f=R^2+t,\qquad
g=R^3+\frac32tR+\frac38,\\
y^2=x^3-3fx+2g,\qquad G=(R,\sqrt3/2),\qquad
\mathcal M=\mathcal O([8G]-[O]),
\end{gathered}
\]
with source translation by \(-G\), selected whole nodal fibre at
\(t=t_c\), and original height \(q_{\rm sel}=t-t_c\). Retain the
original order-one cap at \(t=0\), its positive norm
\(\psi(\widehat\gamma)=1\), every prescribed source and physical
integer in its full period reference, every prescribed integral
lattice-basis selected fan, and an original admitted fixed nonzero
scalar \(\lambda\).
The lifted family is the fixed-$C_*$ continuation of
\cref{night:P3K.116}:
\[
 B_\lambda(p,v_M)=
 \bigl(p-G,\lambda C_*(t-t_c)B_{8,1}(p)v_M\bigr),
 \qquad B_{8,1}=f_2^8\frac{f_8\circ t_{-G}}{f_8}.
\]
The chronological tangent quotients and the single nonzero $C_*$
are fixed by the prescribed initial source and Cartier frames.

On the original ordered selected-source cover
\[
a=-u/3-u^4/9,\qquad t_c=1/(2u)-u^2/3,\qquad
v^2=4u^3+3,
\]
the ordinary-cap collision locus consists of exactly twenty-four
marked points:
\[
\begin{gathered}
16a^3+3=0,\\
2187-16u^3(u^3+3)^3
=-(2u^3-3)(8u^9+84u^6+342u^3+729)=0.
\end{gathered}
\]
Each of the twelve distinct \(u\)-values has two distinct ordered
\(v\)-values. The six points with \(u^3=3/2\) have \(t_c=0\);
the other eighteen have \(t_c\ne0\). These two loci and the other
excluded source boundaries have the disjointness and full local
coordinate inverses proved in OC45.1--2 and ORC.1--2.

Every positive meridian at any of these twenty-four marked points,
closed by its actual original source, line and cap-reference
endpoint, has
\[
[\Phi]=0\qquad\text{in }\Theta_7=\mathbb Z/28.
\]
Every inverse meridian has the same zero class through the
reciprocal source, Cartier, cap, endpoint and transport maps.

The comparison giving this return is an explicit family map.
At each collision let \(s=t_n(a)\) be the actual moving nodal
value, with \(s(a_*)=0\) and \(ds/da=-2a_*\ne0\). Its full
source chart has \(XY=t-s\), the inverse
\[
t=s+XY,\qquad y=(X+h_nY)/2,\qquad
2(x-\xi)\sqrt{x+2\xi}=h_nY-X,
\]
and the complete functions \(\xi,h_n\) of OC45.8--11.
In the actual point coefficient \(V=f_{P^8}(Z;Q)v_M\), the
original periods are
\[
h:(Z,V)\longmapsto(QZ,P^8V),\qquad
B:(Z,V)\longmapsto(P^{-1}Z,CV),\qquad
C=\lambda(t-t_c)\beta .
\]
Here \(Q=(t-s)\kappa\), \(P\), \(\beta\), the fixed Hom
normalization, and all higher coefficients remain the original
functions. At the central node,
\[
P_*=-2\pm\sqrt3,\qquad
\kappa(0,0)=\frac1{288a_*^2},\qquad
\beta(0,0)=C_*\frac9{a_*^2}.
\]
Their full Cartier maps and their target-point inverses occur in
the construction.

Move the reference cap to the protected ordinary point \(t=d\)
by the full based integral and holomorphic comparison
OC45.39--56. Its gluing is
\[
\mathcal F_w(Z,V)=(e^{-w\log P}Z,e^{w\log C}V),\qquad
\Psi_d=\mathcal F_{\log(t-d)/(2\pi i)},\qquad
\Psi_d^{-1}=\mathcal F_{-\log(t-d)/(2\pi i)}.
\]
These maps extend on the original source, point/O modules,
infinity components and every prescribed selected fan by
OC45.29--35. The resulting family \(Y\) is smooth and proper
over the full \(s\)-disc with the original oriented sphere fibre.

On the original cap, the complete retained period is
\[
\gamma_{\rm ref}
=\frac{(-\log P,\log C)}{2\pi i}
+k\,\frac{(\log Q,8\log P)}{2\pi i}
+r\alpha+d_\delta\delta .
\]
Its positive parameter jump is \(k\alpha\) at the eighteen
unselected-cap points and \(k\alpha+\delta\) at the six
selected-cap points. The actual integral preimages, based-cell
inverse and Lie-cochain solution carry each of those full jumps.
With \(F_\theta:X_\theta\to Y_\theta\) the resulting comparison
and \(E_X,E_Y\) the forward marked endpoints, they give
\[
F_1E_XF_0^{-1}=CE_Y,\qquad
C\in\mathbb C^*_{\rm physical}.
\]
For the actual transport \(P_Y:Y_0\to Y_1\), the full return
identity is
\[
F_0\bigl(E_X^{-1}F_1^{-1}P_YF_0\bigr)F_0^{-1}
=E_Y^{-1}C^{-1}P_Y .
\]
The explicit physical isotopy \(e^{\theta\log C}\) and the
disc transport prove the stated class with its prescribed
endpoint. The argument-preserving radial scalar comparison
of OC45.10 carries these same maps to every original admitted
fixed scalar. The complete scope is these twenty-four
ordinary-cap collision meridians and their inverses.
\end{theorem}
\begin{proof}

% Expanded verified cp36 input: notes/part3_os45_cap_collisions_proof.tex
\par\bigskip\noindent\textbf{Complete supporting proof 1: OC45.}\par


% Expanded verified cp36 input: notes/part3_os45_cap_collisions_oc45_proof.tex
\emph{Proof roadmap.}
OC45.1-\/-5 constructs the local enlargement for the eighteen
unselected-cap points and its specified extending period.
OC45.6-\/-11 constructs the full original-reference comparison
and treats all twenty-four points. OCP gives the integral
based-cell maps and their contraction, OCGC gives the global
source/line and endpoint comparison, and ORC records the full
source, proper-family and return argument.

\par

\medskip\noindent\textbf{OS45 ordinary-cap collisions: exact local source and cap maps}

\par\nobreak

OC45.1--5 treats the original cap at \(t=0\) in the
OS45 source and constructs the local enlarged-cap filling
for the eighteen unselected points with its specified
extending period. OC45.6--11 supplies the integral and
analytic comparison for the complete prescribed reference
and all twenty-four points. OCP, OCGC and ORC give the
full lattice, global endpoint and proper-family arguments.

\par

\medskip\noindent\textbf{OC45.1. The twenty-four marked points and their two actual strata}

\par\nobreak

Retain
\[
\begin{gathered}
 R=t^2+a,\quad f=R^2+t,\quad
 g=R^3+\frac32tR+\frac38,\\
 y^2=x^3-3fx+2g,\quad G=(R,\sqrt3/2),
 \qquad q_{\rm sel}=t-t_c.
\end{gathered}
 \tag{OC45.1}
\]
The selected source cover is
\[
 a=-u/3-u^4/9,\quad t_c=1/(2u)-u^2/3,\quad
 v^2=4u^3+3.
 \tag{OC45.2}
\]
The ordinary cap becomes singular exactly at
\[
 16a^3+3=0,\qquad
 2187-16u^3(u^3+3)^3=0.
 \tag{OC45.3}
\]
Writing \(w=u^3\), its complete factorization is
\[
 2187-16w(w+3)^3
       =-(2w-3)(8w^3+84w^2+342w+729).
 \tag{OC45.4}
\]
The twelve distinct \(u\)-values each have two ordered values of \(v\).
At \(u^3=3/2\), the actual selected height satisfies \(t_c=0\).
This gives six selected-cap collision points.
The other three \(w\)-roots give eighteen points with \(t_c\ne0\);
the node at the original cap is then an unselected source node.
The two sets are disjoint by their polynomial gcd. Both are also
disjoint from the selected-cusp, infinity-\(I_9\), and unselected-II
loci of O45.2.

OC45.2--5 constructs the local enlargement for the
eighteen unselected-cap points. The six selected-cap
points retain the full rank-two selected/cap interaction
in OC45.2--4; OC45.6--11 treats both strata through the
protected-cap comparison.

\par

\medskip\noindent\textbf{OC45.2. Original nodal source, moving height and complete point line}

\par\nobreak

Fix one of the eighteen points and write \(a_*\) for its source value.
Then \(a_*^3=-3/16\), the original node is
\((x,y,t)=(-a_*,0,0)\), and the nodal cubic is
\[
 y^2=(x+a_*)^2(x-2a_*).
 \tag{OC45.5}
\]
The original finite discriminant factor is
\(D=48a(t^2+a)^2+8t^3+72at+9\). At this node
\[
 D_t=72a_*,\quad D_a=144a_*^2.
 \tag{OC45.6}
\]
Thus its actual critical base value is a single-valued holomorphic
function \(s=t_n(a)\), with
\[
 s(a_*)=0,\qquad \frac{ds}{da}(a_*)=-2a_*\ne0.
 \tag{OC45.7}
\]
The inverse \(a=a(s)\) is fixed by this original implicit equation.
Set \(q_n=t-s\), keeping the original cap coordinate \(t\).

Choose \(\xi(q_n,s)=\sqrt{f(s+q_n,a(s))}\) with value
\(\xi(0,0)=-a_*\), and put
\[
 h_n(q_n,s)=
  \frac{2(g(s+q_n,a(s))-\xi(q_n,s)^3)}{q_n}.
 \tag{OC45.8}
\]
Its numerator vanishes at \(q_n=0\).
At the collision
\[
 f_t=1,\quad g_t=3a_*/2,\qquad h_n(0,0)=6a_*\ne0.
 \tag{OC45.9}
\]
The actual source chart and inverse are
\[
 X=y-(x-\xi)\sqrt{x+2\xi},\qquad
 Y=h_n^{-1}\bigl(y+(x-\xi)\sqrt{x+2\xi}\bigr),\qquad XY=q_n,
 \tag{OC45.10}
\]
\[
 q_n=XY,\quad t=s+XY,\quad y=(X+h_nY)/2,\quad
 2(x-\xi)\sqrt{x+2\xi}=h_nY-X.
 \tag{OC45.11}
\]
The last implicit derivative is \(2\sqrt{3\xi}\ne0\).
Consequently these are inverse holomorphic source maps over the
joint \(s\)-disc, not merely formulas on its punctured boundary.

On the central nodal normalization put
\[
 x=S^2+2a_*,\quad y=S(S^2+3a_*),\quad
 c^2=-3a_*,\quad Z=(S-c)/(S+c).
 \tag{OC45.12}
\]
The original section has \(S_G=\sqrt3/(4a_*)\).
Its square is \(-a_*\), so \(c/S_G=\pm\sqrt3\), and therefore
\[
 P_*=Z(G)=-2+\sqrt3\quad\hbox{or}\quad -2-\sqrt3
 \tag{OC45.13}
\]
according to the chosen ordered branch.
The two displayed values are reciprocal. Neither is zero,
one, or a root of unity.

The original Miller functions give a useful full nodal identity.
For an arbitrary nonzero node parameter \(c\) and marked
\(P=Z(G)\), the original tangent normalization at \(O\) gives
\[
 f_m(Z)=(-2c)^{m-1}\frac{1-P^m}{(1-P)^m}\,
       \frac{(Z-P)^m}{(Z-P^m)(Z-1)^{m-1}},
       \qquad m=2,4,8.
 \tag{OC45.14}
\]
Indeed the horizontal divisor is the required one, the values
at the two node branches agree, and its original principal
coefficient at \(O\) is \(S^{m-1}\), with
\(S=c(1+Z)/(1-Z)\). Alternatively the three original tangent
quotients successively give the same constants.
The whole chronological coefficient is
\[
 B_{8,1}(Z)=
  \left(-2c\frac{1+P}{1-P}\right)^8
  \frac{(Z-P)(Z-P^8)}{(Z-P^9)(Z-1)}.
 \tag{OC45.15}
\]
For the actual point coordinate
\(V_{\rm pt}=f_{P^8}(Z)v_M\),
where \(f_b=(1-Z/b)/(1-Z)\) on the nodal normalization,
the directed coefficient is
\[
 B_{8,1}(Z)\frac{f_{P^8}(Z/P)}{f_{P^8}(Z)}
       =(-2S_G)^8.
 \tag{OC45.16}
\]
All four point/O factors are retained before this cancellation.
At the cap node this value is
\[
 (-2S_G)^8=\frac{81}{256a_*^8}=\frac9{a_*^2}.
 \tag{OC45.17}
\]
This also verifies the independent selected-node formula:
at \(S_G=-\sqrt3/(2u)\), the value in OC45.16 is \(81/u^8\).

Let \(Q=q_n\kappa(q_n,s)\) and \(P=P(q_n,s)\) be the full
inverse-Tate source functions with this ordered branch and
the original reference scale. Their first coefficient is
\[
 \kappa(0,0)
   =-\frac{h_n(0,0)}{1728\xi(0,0)^3}
   =\frac1{288a_*^2}.
 \tag{OC45.18}
\]
The full physical point-line coefficient is defined, with
its reciprocal, by evaluation of the actual Miller/theta
module ratio at a protected regular point. Write it as
\(\beta(q_n,s)\), including the original normalization
\(C_*\). Then \(\beta(0,0)=C_*9/a_*^2\); its full function
is retained below. Since \(q_{\rm sel}(0,0)=-t_c(a_*)\ne0\),
the complete suspension multiplier
\[
 C(q_n,s)=\lambda q_{\rm sel}\beta(q_n,s)
 \tag{OC45.19}
\]
is an extending unit. The original period arrows are
\[
 h:(Z,V_{\rm pt})\mapsto(QZ,P^8V_{\rm pt}),\qquad
 B:(Z,V_{\rm pt})\mapsto(P^{-1}Z,CV_{\rm pt}).
 \tag{OC45.20}
\]
No central value replaces \(P,\kappa,\beta\), or \(C\) in
these arrows.

\par

\medskip\noindent\textbf{OC45.3. A complete extending local suspension flow}

\par\nobreak

Choose the actual analytic logarithms of the nonzero
functions \(P\) and \(C\) on the small closed joint
\((q_n,s)\)-bidisc. Their values at its reference point
are prescribed. For a complex number \(w\), define
\[
 \mathcal F_w(Z,V_{\rm pt},q_n,s)
   =\left(e^{-w\log P}Z,\ e^{w\log C}V_{\rm pt},q_n,s\right).
 \tag{OC45.21}
\]
Its inverse is \(\mathcal F_{-w}\). The two maps commute
with \(h\) and \(B\), retain the original \(q_n\), and
\[
 \mathcal F_{w+1}=B\mathcal F_w,\qquad
 \mathcal F_{w+w'}=\mathcal F_w\mathcal F_{w'}.
 \tag{OC45.22}
\]
These are identities of the full coefficients.

They extend through the whole nodal source and point
line. On the Tate chain
\(X_i=Z/Q^i,\ Y_i=Q^{i+1}/Z\), the source arrow is
\[
 (X_i,Y_i)\mapsto
       (e^{-w\log P}X_i,\ e^{w\log P}Y_i).
 \tag{OC45.23}
\]
It preserves \(X_iY_i=Q\) and has the written reciprocal
inverse. The exact \(Q=q_n\kappa\) coordinate change
then carries it to the original source chart.
In the original point frame its line coefficient is
\[
 e^{w\log C}\,
       \frac{f_{P^8}(Z)}{f_{P^8}(e^{-w\log P}Z)}.
 \tag{OC45.24}
\]
The defining numerator and denominator modules are
the source and target point/O modules. Their graphs
are carried by the displayed source arrow, so this
is an isomorphism at their zeros and poles as well.
At each node both point sections lie in its smooth
identity component, and the Tate-chain coefficient
in OC45.21 is an actual unit. Thus the complete
map extends across the node, with its inverse.
This proves the local flow as a source/line
isomorphism, not just an equality of a central
coefficient.

In additive logarithmic coordinates its period at
time one is
\[
 \gamma_0=\frac1{2\pi i}
               \binom{-\log P}{\log C}.
 \tag{OC45.25}
\]
The subscript names this particular extending lift.
The prescribed original cap-period lift has not
yet been identified with it.

\par

\medskip\noindent\textbf{OC45.4. Exact enlargement of the cap with this period}

\par\nobreak

For the period \(\gamma_0\), the cap-to-ordinary
map is the original formula, now with its full
extending flow:
\[
 \Psi_0(\zeta,t,s)
       =\mathcal F_{\log t/(2\pi i)}(\zeta,t,s),
 \qquad
 \Psi_0^{-1}
       =\mathcal F_{-\log t/(2\pi i)}.
 \tag{OC45.26}
\]
A logarithm branch change multiplies by \(B\)
by OC45.22. Hence both maps are single-valued
on the actual suspension quotient.
For \(t\ne0\) this includes the moving source
node \(t=s\); the extension follows from
OC45.23-\/-24. Thus a fixed larger cap disc may
contain the moving node.

Here is the precise comparison with the usual
small ordinary cap for \(s\ne0\). On the
small cap use the identity. On the intervening
ordinary region use \(\Psi_0^{-1}\) into the
large cap, and outside the large cap use the
identity on the ordinary region. On the small
cap overlap the original gluing is \(\Psi_0\),
so the two maps compose to the identity.
On the large-cap boundary its gluing is
again \(\Psi_0\), giving the same equality.
These piecewise maps therefore glue to a
holomorphic isomorphism. The reciprocal
piecewise maps give its inverse. Their full
source, point-module and scalar coefficients
are precisely OC45.21-\/-26.

The large cap copy is the original nodal
source/line quotient, with \(q_n=t-s\).
At \((t,s)=(0,0)\) its source total space is
smooth by OC45.10-\/-11; its projection to \(s\)
is a submersion. The lift \(B\) is uniformly
contracting for a common sufficiently small
fixed scalar, since OC45.19 is nonzero and
the line-module map is holomorphic on the
compact local source. Its line-complement
quotient is therefore smooth and proper over
the closed cap disc. Gluing its fixed outer
annulus by OC45.26 gives an actual cap
family over the full \(s\)-disc. All other
source, selected-fan and infinity pieces
retain the O45 maps on this disc.

For reference, the exact base blowup on the
punctured comparison is
\[
 t=s\tau,\qquad
 Q=s(\tau-1)\kappa(s(\tau-1),s),\qquad
 \log t=\log s+\log\tau.
 \tag{OC45.27}
\]
The additive coordinate
\(\bar\zeta=\zeta_{\rm cap}
+(\log s)/(2\pi i)\,\gamma_0\)
then has ordinary transition
\(\zeta_{\rm ord}=\bar\zeta+
(\log\tau)/(2\pi i)\,\gamma_0\).
Both inverse translations retain the complete
period vector. Around the \(s\)-loop this
coordinate changes by precisely \(B\);
it is already the period of the quotient.
No fractional source-power coefficient is
declared to extend as a single-valued
function on the original source chart.

\par

\medskip\noindent\textbf{OC45.5. The exact remaining prescribed-reference equation}

\par\nobreak

On a cut parameter annulus, compare the
prescribed original cap period with OC45.25.
They are lifts of the same original
suspension map on the torus. In the complete
original source/physical lattice their
difference is
\[
 \gamma_{\rm ref}-\gamma_0
        =k\,h+r\,\alpha+d\,\delta,\qquad k,r,d\in\mathbb Z,
 \tag{OC45.28}
\]
where \(h\) is the full source period in
OC45.20, \(\alpha\) is its source circle
period and \(\delta\) the physical circle.
Every one of these integers is retained;
the difference has suspension coordinate
\(\psi=0\). Replacing the cap map by
OC45.26 without the actual comparison
for OC45.28 would change the prescribed
cap datum.

The original five source matrices of K10 give
\(\sum_i\operatorname{im}(T_i^{-t}-I)=\ker\psi\)
as an integral lattice. OC45.6--11 realizes this equality
by explicit based preimages, extending source/point-line
arrows at all five fillings, a parameter-dependent Lie
cochain with its full logarithms, and the directed endpoint
square. Those maps carry the prescribed original cap
reference into the protected-cap family.

\par

\medskip\noindent\textbf{OC45.6. The original extending translation sheaf on all five fillings}

\par\nobreak

The opening status and OC45.1-\/-5 above are retained as the completed
local-stage record. The following construction supplies the original
reference comparison left open there. It also moves the cap to a
protected ordinary point, using explicit maps, and therefore treats
the six selected-cap points as well as the eighteen unselected-cap
points. The cap retains its original order one and positive norm
\(\psi(\widehat\gamma)=1\). No additional root or finite character is
introduced.

Fix the original five source discs and the prescribed selected fan.
All source, point-line and scalar functions in this section are the
ones in O45.1-\/-5, with the full functions of OC45.18-\/-20 on a nodal
disc. A translation on a regular source torus is written in its
actual point coordinates as
\[
 \mathcal T_{A,D}(Z,V,q,s)
       =(e^A Z,e^D V,q,s),\qquad
 \mathcal T_{A,D}^{-1}=\mathcal T_{-A,-D}.
 \tag{OC45.29}
\]
Here \(A,D\) are the two coefficients of the actual Lie vector;
they can depend on the original base and parameter. The scalar
\(\lambda\), the source height and every point-line function remain
fixed in this arrow. In the original meromorphic point section the
line coefficient is exactly
\[
 v_M\longmapsto e^D
       \frac{f_{P^8}(Z;Q)}{f_{P^8}(e^A Z;Q)}v_M.
 \tag{OC45.30}
\]
The inverse evaluates its reciprocal at \(e^A Z\). The numerator and
denominator are the actual source and target point/O modules, so
this is their module isomorphism also at their zeros and poles.
This same formula gives the ordinary torus exponential, using the
complete inverse marked uniformization.

At an unselected node the source-chain coordinates in OC45.23
transform as \((X_i,Y_i)\mapsto(e^A X_i,e^{-A}Y_i)\).
Thus both the flow and its inverse extend on the whole source
chain and its genuine point line. At the selected zero the original
periods have the full form
\[
 h=(\kappa q,P^8),\qquad
 B=(P^{-1},\lambda q\beta)
 \tag{OC45.31}
\]
as multiplicative source/line columns. The two valuation vectors
are \((1,0)\) and \((0,1)\). For every actual character
\(Z^rV^d q^e\) on every cone of the prescribed periodic fan,
OC45.29 multiplies that coordinate by \(e^{rA+dD}\).
These are nonzero analytic units at its unchanged original height,
with inverse \(e^{-rA-dD}\). They agree on all toric overlaps and
commute with the full periods, including \(\kappa,P,\lambda,\beta\).
This proves the extension on each prescribed fan by its actual
coordinate map.

At the original infinity put \(z=1/t\). O45.4 gives the full
\[
 P_\infty=z^3\pi(z,s),\qquad Q_\infty=z^8\kappa_\infty(z,s),
 \qquad b_0=P_\infty^8/Q_\infty^3,
 \qquad b_0(0,s)=-1/(9a(s)^3).
 \tag{OC45.32}
\]
All displayed base coefficients here remain analytic units near
the cap-collision disc. The exact product identity
\[
 f_{Q_\infty^3b_0}(Z;Q_\infty)
       =(-Z/b_0)^3Q_\infty^{-6}f_{b_0}(Z;Q_\infty)
 \tag{OC45.33}
\]
identifies the original point coordinate with the multidegree-zero
coordinate \(V_0=f_{b_0}(Z;Q_\infty)v_M\). Consequently OC45.29
acts there by
\[
 Z\longmapsto e^A Z,\qquad V_0\longmapsto e^{D-3A}V_0.
 \tag{OC45.34}
\]
The factors \(b_0^3Q_\infty^6\) in the reciprocal coordinate
comparison are retained; they cancel in this displayed action
because the original base is fixed. In that frame the actual
period reindex is
\[
 H=h^{-1}B^{-3},\quad K=h^3B^8,\qquad
 h=H^8K^3,\quad B=H^{-3}K^{-1}.
 \tag{OC45.35}
\]
It has determinant one, source valuations one and zero, and the
complete line multiplier of \(K\) is the original contracting
unit. Formula OC45.34 commutes with both full columns. On the
primitive rank-one chain it and its inverse multiply every
coordinate by its actual unit character. Thus it extends through
the entire infinity filling. No constant term is substituted for
either period or the point-to-tensor factor.

Let \(\mathscr E\) denote this holomorphic Lie bundle on the
original base, and \(\mathscr H\) the kernel of its exponential
in the sheaf \(\mathscr A\) of the extending translations just
constructed. The period sheaf is \(j_*\Lambda\), with local stalk
\(\ker(H_i-I)\). Here is its actual local identification. The
logarithmic radial coefficients of the full periods have exactly
the primitive valuation maps of OC45.31 and OC45.35. A period
extends in the Lie bundle precisely when these radial coefficients
vanish. The remaining angular periods are the logarithms of
the extending unit actions in OC45.29, with the full point frame
OC45.30. This is exactly the kernel of the original monodromy.
On two slit sectors and the centre the marked circle complex
\(\Lambda\xrightarrow{H_i-I}\Lambda\) has its degree-one
coinvariant term cancelled by the identical annular term.
Thus the invariant extension has no higher cohomology on a
source disc. This is an integral cancellation in the stated
circle coordinates.

Logarithms of the two nonzero unit multipliers give the local
inverse to the exponential. The primitive period reindexes above
remove the integral radial orders before these logarithms are
taken. The resulting exact sequence is therefore the sequence
of the actual source and point-line maps:
\[
 0\longrightarrow\mathscr H\longrightarrow\mathscr E
       \xrightarrow{\operatorname{Exp}}\mathscr A\longrightarrow0.
 \tag{OC45.36}
\]
Only the existing source fillings and their same fan occur here.

The Lie source line can be calculated directly on OS45. In the
original minimal infinity chart \(x=z^{-2}X\), \(y=z^{-3}Y\),
the relative invariant differentials satisfy
\[
 \omega_t=\frac{dx}{2y}
           =z\frac{dX}{2Y}=z\omega_\infty .
 \tag{OC45.37}
\]
Their dual vectors therefore have
\(e_\infty=t^{-1}e_0\). At a nodal point the relative dualizing
differential is a unit times \(dX_i/X_i=-dY_i/Y_i\), and its
dual is the extending field \(X_i\partial_{X_i}-Y_i\partial_{Y_i}\).
At the selected fan this is the source unit action already
displayed; at infinity it is carried by OC45.33-\/-35. Hence the
source Lie line is literally \(\mathcal O(-1)\) with these frames.
The physical line is the global field \(\delta\) whose flow is
multiplication by \(e^{2\pi iz}\). The complete Lie extension is
\[
 0\longrightarrow\mathcal O\delta\longrightarrow\mathscr E
       \longrightarrow\mathcal O(-1)\longrightarrow0.
 \tag{OC45.38}
\]
In point-module charts its off-diagonal entry is the source
invariant vector applied to the logarithm of the entire actual
line transition, including OC45.30 and OC45.33, divided by the
chosen logarithmic normalization. No central line constant
replaces this derivative.

\par

\medskip\noindent\textbf{OC45.7. The complete integral comparison in the actual based cells}

\par\nobreak

Use exactly K10\textquotesingle s five matrices \(T_i\), in order
\((\psi,u,w,\delta)\), and set
\[
 H_i=T_i^{-t},\quad N_i=H_i-I,\quad
 P_0=I,\quad P_i=H_0\cdots H_{i-1},\quad
 M_i=P_iN_i=P_{i+1}-P_i.
 \tag{OC45.39}
\]
The product is \(H_0\cdots H_4=I\). The complete source basis
and transported cells are those of K10, not an untransported
sum of abstract image lattices. For selected zero at site
\(\sigma\in\{1,2,3,4\}\) put
\[
 d_1=e_\psi,\quad d_2=e_\psi+2e_w,\quad
 d_3=e_\psi+e_w,\quad d_4=e_\psi .
 \tag{OC45.40}
\]
Direct inversion of each of the four actual central-entry choices
gives
\[
 N_\sigma d_\sigma=e_\delta,\qquad
 N_1(-e_u)=e_w,\qquad N_3e_w=e_u+8e_\delta.
 \tag{OC45.41}
\]
For an arbitrary retained difference
\(v=Ae_u+Be_w+Ce_\delta\), start with all \(r_i=0\), set
\[
 r_1=-Be_u,\qquad r_3=Ae_w,\qquad
 r_\sigma\mathrel{+}=(C-8A)d_\sigma .
 \tag{OC45.42}
\]
Repeated indices receive the indicated sum. Then
\(\sum_iN_ir_i=v\), with no division. Its actual based-cell
preimages and inverse are
\[
 t_j=r_j-\sum_{i>j}N_i r_i,\qquad
 r_j=t_j+\sum_{i>j}N_i r_i\quad(j=4,3,2,1,0).
 \tag{OC45.43}
\]
The identity \(\sum_{j<i}M_j=P_i-I\) proves
\(\sum_jM_jt_j=\sum_iN_ir_i=v\).
This is an upper triangular integral coordinate isomorphism
with identity diagonal.

For clarity the full cone sign is retained too. With exterior
one-cochains \(p_0,\ldots,p_3\) and annular zero-cochains \(z_i\),
the degree-two boundary is
\[
 b_i=p_i-N_iz_i\ (i<4),\qquad
 b_4=-H_4\sum_{i<4}P_ip_i-N_4z_4.
 \tag{OC45.44}
\]
The coordinate \(v_*=\sum_iP_i b_i\) cancels every exterior
\(p_i\), and its boundary is \(-\sum_iM_i z_i\).
Thus the positive difference \(v\) in OC45.42 is produced
by exactly \(z_i=-t_i\). These are the restrictions along
the actual based cuts. Smooth cap discs add their transported
period vector to \(v_*\); their monodromy is the identity.

We also retain the entire degree-one calculation. The primitive
local image basis is
\[
 g_0=e_w+3e_\delta,\quad g_1=e_w,\quad
 g_2=e_u+3e_w+16e_\delta,\quad
 g_3=e_u+8e_\delta,\quad g_4=e_u-3e_w,
 \tag{OC45.45}
\]
with one additional \(e_\delta\) at the selected site.
Record these six coefficients in that order, placing the
additional coefficient last. In these actual integral bases
the parabolic complex has matrices
\[
 D=\begin{pmatrix}
 -3&-8&0&0\\0&-1&0&0\\-2&-3&1&0\\
 -1&0&1&0\\0&3&1&0\\1&0&0&0
 \end{pmatrix},\qquad
 B=\begin{pmatrix}
 0&0&0&0&0&0\\
 0&0&1&-2&1&0\\
 1&1&-6&9&-3&0\\
 3&0&-8&8&0&1
 \end{pmatrix}.
 \tag{OC45.46}
\]
Here \(D(v)=(N_iv)_i\); the first five columns of \(B\) are
\(P_i g_i\), and its last is \(P_\sigma e_\delta=e_\delta\).
These matrices are independently obtained for every selected
site after the displayed placement of the extra generator.
For every integral cocycle,
\[
 Bz=0\quad\Longrightarrow\quad
 z=D(z_5,-z_1,z_3+z_5,0)^t.
 \tag{OC45.47}
\]
Let \(D'\) be the first three columns of \(D\), \(B'\) the
last three rows of \(B\), and define
\[
 Lz=(z_5,-z_1,z_3+z_5)^t,\qquad
 J=\begin{pmatrix}
 -24&-8&3\\0&0&0\\-9&-3&1\\0&0&0\\
 10&3&-1\\0&0&0
 \end{pmatrix}.
 \tag{OC45.48}
\]
Their full integer identities are
\[
 LD'=I_3,\quad B'J=I_3,\quad D'L+JB'=I_6,\quad B'D'=0.
 \tag{OC45.49}
\]
They give explicit chain contraction and splitting maps,
including all torsion information. With the actual local
invariant-stalk cancellation of OC45.36, they prove
\[
 H^0(\mathscr H)=\mathbb Z\delta,\qquad
 H^1(\mathscr H)=0,\qquad
 H^2(\mathscr H)\xrightarrow{\ \psi\ }\mathbb Z.
 \tag{OC45.50}
\]
The complete independent derivation, including the twenty
original inverse matrices and the full twenty-by-twenty map
OC45.43, is retained in OCP.1-\/-6 and its executable certificate.

\par

\medskip\noindent\textbf{OC45.8. The holomorphic comparison with every original reference integer}

\par\nobreak

Compare two actual order-one cap constructions on this same
source, line, scalar and selected fan. Each cap has the original
positive norm. The first is the prescribed cap at \(t=0\).
The second will be placed at a protected ordinary value \(t=d\)
in OC45.10. On a cut parameter disc join their based cap collars
by a specified path in the original regular source base.
Transport the complete original lattice along this path.
The difference of their two cap cocycles in the coordinate
of OC45.44 is an integral vector \(v\) with \(\psi(v)=1-1=0\).
Every other component of \(v\) is retained as \(A,B,C\) in
OC45.42. The two base logarithms are \(\log t\) and
\(\log(t-d)\), respectively; neither is replaced by the other.
The actual collar maps are their full plus-log translations.

Take logarithms of the actual difference of these translation
gluing maps on slit overlaps. Their triple-overlap discrepancy
is their full period matrix \(\Pi(t,s)\) times the integral
two-cocycle just described. Subtract the periods of the precise
total one-cochain OC45.42-\/-44. The actual \(e_\psi\) terms
in \(d_\sigma\) are the original suspension deck, with its
whole source and line multiplier OC45.31 or OC45.35.
They are inserted at both ends of the corresponding overlap.
The corrected Lie cocycle is single valued on each annulus:
its two slit branches differ by the cancelled \(N_i z_i\).
The full source periods in \(\Pi\), including \(\log Q\),
\(\log P\), the physical period and the suspension logarithm,
remain in these subtractions. No period is declared constant
as a parameter function.

Here is a reproducible analytic solution of this remaining
Lie cocycle. Use the actual two based Lie frames of OC45.38.
Their transition has the form
\[
 \mathcal L(t,s)=
     \begin{pmatrix}t^{-1}&0\\a(t,s)&1\end{pmatrix}.
 \tag{OC45.51}
\]
The entire coefficient \(a(t,s)\) is the point-module derivative
specified after OC45.38. Holomorphic lifts of the source frame
on the two affine base charts are obtained by the scalar Cauchy
operator on those charts applied to their point-frame differences;
this retains the same coefficient on the overlap.
For a holomorphic overlap vector \((f,g)\), expand in Laurent
series on a fixed protected circle \(|t|=\rho\). Its coefficients
are the actual integrals
\[
 f_j(s)=\frac1{2\pi i}\int_{|t|=\rho}f(t,s)t^{-j-1}\,dt .
 \tag{OC45.52}
\]
The equation
\(\mathcal L(u_\infty,v_\infty)^t-(u_0,v_0)^t=(f,g)^t\)
has the explicit solution
\[
 u_0=-\sum_{j\ge0}f_jt^j,\qquad
 u_\infty=\sum_{j<0}f_jt^{j+1},\qquad
 G=g-a(t,s)u_\infty,
 \tag{OC45.53}
\]
\[
 v_0=-\sum_{j\ge0}G_jt^j,\qquad
 v_\infty=\sum_{j<0}G_jt^j.
 \tag{OC45.54}
\]
The negative powers in the infinity components are holomorphic
in \(1/t\), including the \(j=-1\) constant in \(u_\infty\).
Direct substitution proves both equations and keeps the whole
off-diagonal term.

For the original many-disc cover, choose a fixed smooth partition
of unity \(\rho_k\) on its refinement. If the corrected cocycle is
\(\theta_{ij}\) and its Lie transition is \(\mathcal L_{ij}\),
the smooth cochain
\(s_i^{\rm sm}=\sum_k\rho_k\theta_{ik}\), in frame \(i\),
satisfies \(s_i^{\rm sm}-\mathcal L_{ij}s_j^{\rm sm}=\theta_{ij}\).
Its \(\bar\partial\) is a global \(\mathscr E\)-valued form.
On the two affine charts take its componentwise Cauchy primitives,
using fixed compactly supported extensions and the kernel
\((\pi(t-w))^{-1}\). Their difference on the overlap is holomorphic.
Apply OC45.53-\/-54 to this difference to obtain matching primitives.
Subtracting the resulting global primitive from \(s_i^{\rm sm}\)
produces holomorphic \(s_i\) with exactly the same cocycle equation.
All partitions and integration kernels are fixed in the parameter;
every resulting coefficient is holomorphic in \(s\) on the
chosen cut. This is also an explicit proof of \(H^1(\mathscr E)=0\).

Write the original and reference overlaps in the directed convention
\[
 g^X_{ij}=t_{\theta_{ij}}g^Y_{ij},\qquad
 \theta_{ij}=s_i-\mathcal L_{ij}s_j.
 \tag{OC45.55}
\]
Then the full family map and inverse are
\[
 F_i=\operatorname{Exp}(-s_i),\qquad F_i^{-1}=\operatorname{Exp}(s_i),
 \qquad F_i g^X_{ij}=g^Y_{ij}F_j.
 \tag{OC45.56}
\]
The last identity is an equality of the whole original source
and point-line coefficients. Formulae OC45.29-\/-35 carry it to
every ordinary, nodal, infinity and prescribed selected-fan
chart, with reciprocal coefficients in the inverse. On a
smooth cap it is the actual torus exponential before its
plus-log gluing. It therefore includes the cap circle and
the complete annular reference, rather than just the integral
class of that reference.

\par

\medskip\noindent\textbf{OC45.9. The full parameter closing square}

\par\nobreak

The parameter monodromy of the prescribed original period is
visible in the local coordinates; it is not suppressed by
OC45.56. At the original cap \(t=0\) put \(q_n=-s\). On a cut
disc write its complete period as
\[
 \gamma_{\rm ref}=\gamma_0+k h+r\alpha+d_\delta\delta,
 \qquad \gamma_0=(-\log P,\log C)/(2\pi i),
 \tag{OC45.57}
\]
where \(h=(\log Q,8\log P)/(2\pi i)\) and \(\alpha,\delta\)
are the original source and physical circle periods.
All integers are constant on this cut in the actual transported
marking. A positive \(s\)-turn has \(h\mapsto h+\alpha\).
At an unselected-cap collision \(C(0,s)\) is an extending unit,
so \(\gamma_0\) returns; at a selected-cap collision it is
exactly \(-\lambda s\beta(-s,s)\), so
\(\gamma_0\mapsto\gamma_0+\delta\). Therefore the full jumps are
\[
\begin{gathered}
 \Delta\gamma_{\rm ref}=k\alpha
       \quad\hbox{at each of the eighteen unselected-cap points},\\
 \Delta\gamma_{\rm ref}=k\alpha+\delta
       \quad\hbox{at each of the six selected-cap points}.
\end{gathered}
 \tag{OC45.58}
\]
The analytic logarithms of \(P,\kappa,\beta\) and every other unit
have their actual continued values; their integer changes are
already included in the chosen \(k,r,d_\delta\) marking.
After its full based transport, each vector in OC45.58 is
inserted into OC45.42-\/-44. It has zero \(\psi\)-coordinate.
This explicitly carries the parameter endpoint of the original
cap, including the additional physical period at the six points.

The group of differences of the local parameter comparisons
can be calculated from the same exact maps. A global section
of OC45.38 projects to a global section of \(\mathcal O(-1)\).
In the two charts its equation is \(u_0=t^{-1}u_\infty\);
the regular Laurent powers on the two sides have no common
nonzero term. Hence it is zero. The remaining component is a
constant physical vector. By OC45.36 and OC45.50,
\[
 0\longrightarrow\mathbb Z\delta\longrightarrow\mathbb C\delta
 \longrightarrow H^0(\mathscr A)\longrightarrow0,\qquad
 [z\delta]\longmapsto(v_M\mapsto e^{2\pi iz}v_M).
 \tag{OC45.59}
\]
The last arrow is onto because the actual integral \(H^1\)
is zero by OC45.47-\/-49. Its kernel is exactly the primitive
physical period. Thus every ratio of the constructed global
comparisons is the indicated physical \(\mathbb C^*\) action,
including its literal action and inverse on every cap and fan.

For completeness, these local comparisons have an explicit
closed holomorphic choice on a thin parameter annulus.
Write \(F_\alpha=c_{\alpha\beta}F_\beta\) on parameter overlaps.
The coefficients \(c_{\alpha\beta}\) are the physical scalars
of OC45.59, retaining the original endpoint substitutions.
For a partition of unity \(\rho_\beta\) put
\[
 \mathfrak a_\alpha=\sum_\beta\rho_\beta\,d\log c_{\alpha\beta},
 \qquad
 \mathfrak a_\alpha-\mathfrak a_\gamma=d\log c_{\alpha\gamma}.
 \tag{OC45.60}
\]
Angular parallel transport gives a nonzero comparison with
holonomy \(c\in\mathbb C^*\). Choose a logarithm of this one
number and a smooth function \(\chi(\theta+1)=\chi(\theta)+1\)
which is constant on the endpoint collars. Multiplication by
\(\exp[-\chi(\theta)\log c]\) makes it periodic with matching
derivatives. Its extension to a thin radial annulus is the
actual connection equation
\[
 \partial_\rho z_\alpha
       =-\mathfrak a_\alpha(\partial_\rho)z_\alpha,\qquad
 z_\alpha(\rho,\theta)=z_\alpha(\rho_0,\theta)
       \exp\left[-\int_{\rho_0}^{\rho}
                 \mathfrak a_\alpha(\partial_r)\,dr\right].
 \tag{OC45.61}
\]
Here the comparison is \(z_\alpha F_\alpha\). The transition
rule in OC45.60 proves that these solutions agree. Its
antiholomorphic derivative is a single scalar form
\(\bar\partial F^{\rm sm}=\eta F^{\rm sm}\). Extend its
coefficient smoothly with compact support in the parameter
plane and set
\[
 u(s)=\frac1\pi\int_{\mathbb C}
                    \frac{\eta_{\bar w}(w)}{s-w}\,dA(w).
 \tag{OC45.62}
\]
Then \(\bar\partial u=\eta\) on the smaller annulus, and
\(e^{-u}F^{\rm sm}\) is an actual nonzero holomorphic family
comparison. This multiplies only by the already proved
physical action, with inverse multiplication. It does not
discard a source period, a cap integer or a parameter seam.

Equivalently, the endpoint can be retained directly on a cut.
Let \(F_\theta:X_\theta\to Y_\theta\) be OC45.56 with all
the continued periods in OC45.58. Let \(E_X:X_0\to X_1\)
and \(E_Y:Y_0\to Y_1\) be the forward marked endpoints.
The constructed cap endpoint is \(F_1^{-1}E_YF_0\); its
local coefficients and reciprocal are exactly OC45.56,
with the integer correction OC45.58 at the final cut.
The supplied reference endpoint is carried with those same
period decks. Its difference is therefore the actual scalar
\[
 C=F_1E_XF_0^{-1}E_Y^{-1}\in\mathbb C^*_{\rm physical},
 \qquad F_1E_XF_0^{-1}=CE_Y .
 \tag{OC45.63}
\]
This scalar is retained, and its explicit isotopy
\(e^{\theta\log C}\) extends on every original line, fan and cap.
For a smooth parallel transport \(P_Y:Y_0\to Y_1\), choose
the transport on \(X\) through \(F_\theta\). Its inverse-endpoint
return satisfies the exact conjugacy equation
\[
 F_0\bigl(E_X^{-1}F_1^{-1}P_YF_0\bigr)F_0^{-1}
           =E_Y^{-1}C^{-1}P_Y .
 \tag{OC45.64}
\]
Thus the closed-annulus construction and the cut construction
give the same original marked return, with its full endpoint.

\par

\medskip\noindent\textbf{OC45.10. A protected ordinary cap across all twenty-four points}

\par\nobreak

Fix any one of the twenty-four marked points OC45.3. The
source critical fibre at the original cap is the ordinary
node OC45.5, with the exact moving height \(s\) of OC45.7.
This holds at the six points too: there the selected nodal
value is precisely \(t_c=s\). The two branch values in
OC45.13 are nonzero and distinct. Choose constants
\(0<\epsilon\ll |d|\) inside the same source-coordinate
neighbourhood, and a closed disc about \(t=d\) of radius
less than \(|d|/4\). Shrink \(|s|<\epsilon\) so that this
disc is disjoint from every source critical value and from
the selected value. Its original coordinate is \(t-d\);
its source functions are still the full functions of
OC45.1 with \(a=a(s)\).

On this protected disc, \(Q=(t-s)\kappa(t-s,s)\),
\(P=P(t-s,s)\), and \(C=\lambda(t-t_c(s))\beta(t-s,s)\)
are nonzero. Choose their original analytic logarithms
on this joint simply connected disc, fixed at one reference
point. The source/line flow is precisely OC45.21 there.
The cap-to-ordinary gluing and its inverse are
\[
 \Psi_d=\mathcal F_{\log(t-d)/(2\pi i)},\qquad
 \Psi_d^{-1}=\mathcal F_{-\log(t-d)/(2\pi i)}.
 \tag{OC45.65}
\]
A logarithm branch changes this by the entire original
suspension \(B\). The maps therefore descend literally
to the compact torus quotient. Its cap period has norm
one and all its source/physical components are retained
by the displayed logarithms.

Call the resulting family \(Y\). At the original point
\(t=0\) it has no extra logarithmic clutch: it uses the
actual unselected source quotient for the eighteen points,
and the prescribed selected periodic fan for the other
six points. Their charts extend over the parameter centre.
For the unselected source this is OC45.10-\/-11 and
OC45.29-\/-30. For the selected zero the exact full
\(q=t-s\) columns OC45.31 have analytic nonzero
\(\kappa,P,\beta\), with \(P_*\) given in OC45.13.
Each fan chart and its reciprocal unit transition hence
extend with the unchanged \(q\). At original infinity
\(a_*^3=-3/16\), so OC45.32 gives
\(b_0(0,0)=16/27\), a nonzero nonidentity point. The full
source/line maps OC45.32-\/-35 extend there as well.
Every other source critical point is simple and distinct
on this parameter disc, by the exact polynomial gcds
in OC45.1 and the nonzero original discriminant derivative.
The remaining O45 source and Cartier-module charts
therefore extend on the same disc.

This is a smooth proper family, as follows from those
actual local maps. The source total space near the moving
node is \(XY=t-s\), so its coordinates \((X,Y,s)\)
exhibit a submersion to \(s\). At the selected fan the
smooth unimodular cone coordinates together with \(s\)
give the same submersion. The original infinity and
ordinary source graphs retain their smooth local
coordinates and extending analytic units. The cap at
\(d\) is the actual smooth torus disc before OC45.65.
Choose a common strictly contracting scalar for these
compact source pieces. Their contracting fundamental
annuli are compact, and OC45.35 gives the primitive
rank-one period coordinates at infinity. The selected
valuation basis OC45.31 and a finite set of cones
modulo its two periods give a compact fundamental
set there; all remaining multipliers are bounded
nonzero units on the closed parameter disc.
The finitely many indicated compact pieces are glued
along the displayed collars and their inverses.
They give a compact Hausdorff total inverse image of
the closed parameter disc, and the local submersions
agree. This proves properness and smoothness of \(Y\)
without an unlisted piece at the cap collision.

The original cap family \(X\) over the punctured parameter
disc and \(Y\) have the same degree-two class by
OC45.42-\/-50. Their full comparison is OC45.51-\/-65;
in particular the source/physical periods of the original
cap are carried by an actual map and inverse. On a chosen
reference parameter, the member \(X\) is the original
K10 sphere transported through the O45 source, line,
scalar, cap and selected-fan charts. For every prescribed
fan its reference sphere is obtained by the finite
canonical word and the actual common-blowup/two-ruling
maps, primitive Gysin and van Kampen maps of F35,
as retained in O45P.6. Thus \(Y\) has a sphere member.
Its proper smooth disc family transports that member,
with its complex orientation, to every fibre.

The comparison is made at the original fixed nonzero
scalar. The usual contraction domain is a punctured disc.
For a larger admitted scalar on a compact source loop,
the explicit radial path
\(\lambda_\tau=[(1-\tau)+\tau\varepsilon]\lambda_0\),
with \(0<\varepsilon<1\), remains contracting because
the entire suspension multiplier is linear in \(\lambda\).
Apply the same source, point-line and cap maps throughout
this compact interval. It joins the given loop to one
in the common disc domain above, with the same endpoint
periods and no scalar turn. Its smooth cylinder transport
conjugates the two marked returns. This carries the
disc-filling conclusion to that original admitted scalar.

\par

\medskip\noindent\textbf{OC45.11. The marked return and the exact completed scope}

\par\nobreak

For \(Y\), the marked endpoint around the parameter disc
is its identity, because its full source, period, line,
cap and fan maps were constructed over that disc. Smooth
parallel transport of this proper submersion gives a
return isotopic to the identity. The comparison
OC45.63-\/-64 carries the original cap reference and
its actual endpoint to this disc family. The remaining
factor \(C^{-1}\) has the literal global physical
isotopy in OC45.63. Hence each of the twenty-four
ordinary-cap collision meridians in the ordered OS45
source cover has
\[
 [\Phi]=0\quad\hbox{in }\Theta_7=\mathbb Z/28
 \tag{OC45.66}
\]
for the original order-one positive-norm cap, its
complete prescribed period reference and every
prescribed lattice-basis selected fan. The inverse
meridian uses the reciprocal maps and has the same
zero class. The statement concerns this original
source, fixed scalar and cap-reference closing;
the full source and linear markings in the comparison
have been kept throughout.

The local enlargement OC45.21-\/-27 is retained as a
second explicit filling for the eighteen unselected
points and its particular extending period. The
protected-point comparison proves the original
reference assertion and the additional six points;
it does not identify the two periods by notation.
The integral contraction, full Lie-cochain solution,
source/line extension and directed endpoint square above
establish the asserted original-reference return.
The remaining supporting proofs give the complete
based-cell, global-comparison and proper-family calculations.

% End expanded verified cp36 input


\par\bigskip\noindent\textbf{Complete supporting proof 2: OCP.}\par


% Expanded verified cp36 input: notes/part3_os45_cap_collisions_ocp_proof.tex
\par

\medskip\noindent\textbf{OS45 ordinary-cap period lattice: exact based-cell comparison}

\par\nobreak

Use the five original K10 matrices in the retained order
\((\psi,u,w,\delta)\). The following integral preimages
and directed analytic cochain map retain the full
parameter endpoint of the prescribed cap reference.

\par

\medskip\noindent\textbf{OCP.1. Primitive integral images in the original marking}

\par\nobreak

Let \(T_i\), \(0\leq i\leq4\), be exactly K10\textquotesingle s displayed cohomology
matrices and put \(H_i=T_i^{-t}\), \(N_i=H_i-I\). The selected zero is
first at site4. Direct inversion gives
\[
\begin{aligned}
N_0v&=(-3v_\psi-8v_u)(e_w+3e_\delta),\\
N_1v&=-v_u e_w,\\
N_2v&=(-2v_\psi-3v_u+v_w)(e_u+3e_w+16e_\delta),\\
N_3v&=(-v_\psi+v_w)(e_u+8e_\delta),\\
N_4v&=(3v_u+v_w)(e_u-3e_w)+v_\psi e_\delta.
\end{aligned}
\tag{OCP.1}
\]
Every displayed rank-one coefficient row has gcd1. At site4 the two
coefficient rows separately take the values \((1,0)\) and \((0,1)\).
These are the primitive image lattices of the actual local fillings.

For selected site \(s\in\{1,2,3,4\}\), the source K10 rule adds one to
the central entry of \(T_4\) and subtracts one at \(T_s\). Set
\[
d_1=e_\psi,\quad d_2=e_\psi+2e_w,\quad
d_3=e_\psi+e_w,\quad d_4=e_\psi.
\tag{OCP.2}
\]
For every one of these four actual choices the full inverse-transpose
matrices satisfy
\[
N_s d_s=e_\delta,\qquad N_1(-e_u)=e_w,
\qquad N_3e_w=e_u+8e_\delta.
\tag{OCP.3}
\]
Thus for the retained difference
\(v=Ae_u+Be_w+Ce_\delta\), define the five source cochains by
\[
r_1=-Be_u,\qquad r_3=Ae_w,\qquad
r_s\mathrel{+}=(C-8A)d_s,
\tag{OCP.4}
\]
starting with all other entries zero and adding terms at repeated
indices. Their full image is exactly \(\sum_iN_ir_i=v\).
The first row of every \(H_i\) is \((1,0,0,0)\). Consequently
\[
\sum_i\operatorname{im}N_i=\ker\psi,\qquad
\Lambda/\sum_i\operatorname{im}N_i\xrightarrow{\ \psi\ }\mathbb Z
\tag{OCP.5}
\]
is an integral isomorphism, with section \(1\mapsto[e_\psi]\).
All source-cycle and physical-circle coefficients in \(v\) are kept.

\par

\medskip\noindent\textbf{OCP.2. The actual based polygon and its inverse coordinate change}

\par\nobreak

Use the oriented based boundary word \(H_0H_1H_2H_3H_4=I\) of K10.
Write
\[
P_0=I,\qquad P_i=H_0\cdots H_{i-1},\qquad
M_i=P_i(H_i-I)=P_{i+1}-P_i.
\tag{OCP.6}
\]
These are the transported cellular boundary columns. The following
explicit cochain map carries the raw preimages into these columns:
\[
t_j=r_j-\sum_{i>j}N_i r_i,\qquad 0\leq j\leq4.
\tag{OCP.7}
\]
Indeed \(P_i-I=\sum_{j<i}M_j\), and therefore
\[
\sum_jM_jt_j
=\sum_i\{P_iN_i-(P_i-I)N_i\}r_i
=\sum_iN_ir_i=v.
\tag{OCP.8}
\]
The map in OCP.7 is upper triangular over \(\mathbb Z\), with identity
diagonal blocks. Its actual inverse is the descending recursion
\[
r_j=t_j+\sum_{i>j}N_ir_i,\qquad j=4,3,2,1,0.
\tag{OCP.9}
\]
Here is the full cone direction. On the exterior take four one-cochains
\(p_0,\ldots,p_3\); the last restriction is
\(-H_4\sum_{i<4}P_ip_i\). With local annular zero-cochains \(z_i\),
the degree-two boundaries are
\[
b_i=p_i-N_i z_i\quad(i<4),\qquad
b_4=-H_4\sum_{i<4}P_ip_i-N_4z_4.
\tag{OCP.9a}
\]
The integral coordinate \(v_* = \sum_iP_i b_i\) eliminates every
\(p_i\) and sends this boundary to \(-\sum_iM_i z_i\). To produce
the positive difference \(v\) of OCP.8 use exactly \(z_i=-t_i\).
These are the actual exterior/local restrictions and their signs.

Thus the transported cell quotient is literally \(\mathbb Z\) by
\(\psi\). K10\textquotesingle s primitive local specialization maps and the exact
slit-sector invariant-stalk contraction of K78, OCS.1, identify this
cone with the no-finite-root extending period sheaf. Its
degree-two obstruction group is \(H^2(\mathbf P^1,\mathscr H)=\mathbb Z\)
with the specified generator. The ordinary cap vector has
\(\psi(\gamma)=1\), so its actual clutch represents the displayed
generator. Two norm-one cap vectors have the same degree-two class;
OCP.7 and the sign in OCP.9a supply the full integral cochain for their
difference.

\par

\medskip\noindent\textbf{OCP.3. Directed local logarithms and the full line transition}

\par\nobreak

For OC45\textquotesingle s local flow retain its full period
\[
\gamma_0=( -\log P,\log C)/(2\pi i),\qquad
C=\lambda q_{\rm sel}\beta.
\tag{OCP.10}
\]
Express the prescribed difference \(\gamma_{\rm ref}-\gamma_0\)
in the original K10 common period marking before using OCP.4.
The integer \(\psi\)-component is zero. OCP.4-\/-9 subtract precisely
that integral total coboundary from the actual atlas logarithms.
An occurrence of \(e_\psi\) in \(d_s\) is the actual suspension
deck, including its original source and line multiplier. Its full
deck map is retained at both ends of the affected overlap.

After this integral adjustment the difference of the actual annular
translation cocycles has a vector-sheaf lift. The source translation
line is the original \(\mathcal O(-1)\), and the vertical physical
line is \(\mathcal O\). These are the invariant differential and
physical-circle frames in K41/K44; their exponentials are the extending
translations on the primitive source and selected Cox charts. The
off-diagonal entry is retained from differentiating the full original
point-line transition. In two based source charts its coefficient
matrix has the form
\[
L(t)=\begin{pmatrix}t^{-1}&0\\a(t)&1\end{pmatrix}.
\tag{OCP.11}
\]
Here \(a(t)\) is that actual holomorphic overlap coefficient; no value
is substituted for it. For a vector cocycle \((f,g)\) the directed
equation \(L(t)(u_\infty,v_\infty)^t-(u_0,v_0)^t=(f,g)^t\) is solved by
the following fixed-contour Laurent operators:
\[
u_0=-\sum_{n\geq0}f_nt^n,\qquad
u_\infty=\sum_{n<0}f_nt^{n+1},\qquad
G=g-a(t)u_\infty,
\tag{OCP.12}
\]
\[
v_0=-\sum_{n\geq0}G_nt^n,\qquad
v_\infty=\sum_{n<0}G_nt^n.
\tag{OCP.13}
\]
The equations follow by direct substitution, including the entire
\(a(t)u_\infty\) term. The Laurent coefficients are the Cauchy
integrals on the same protected base circle. Thus the solution retains
the full line module. On the protected parameter sector the actual
fixed-contour Cauchy integrals are holomorphic parameter functions.

With gluing convention \(g_{0i}^{\rm ref}=g_{0i}^0\exp(c_{0i})\),
the solved cochain is \(c_{0i}=b_0-b_i\). The directed comparison is
\[
\Psi_i=\exp(-b_i),\qquad
\Psi_0g_{0i}^{\rm ref}=g_{0i}^0\Psi_i,
\qquad \Psi_i^{-1}=\exp(b_i).
\tag{OCP.14}
\]
K44/K78\textquotesingle s actual translation and point-line maps realize this exponential
on the compactified charts; the integral deck terms in OCP.7 are inserted
before exponentiation. This is the full source/line cochain arrow.

\par

\medskip\noindent\textbf{OCP.4. The retained parameter endpoint}

\par\nobreak

OCP.14 preserves the endpoint carried by the supplied logarithms. On a
parameter interval its endpoint comparison is the actual composition
\(\Psi_i(1)\Psi_i(0)^{-1}\), with the supplied original closing maps
on the two sides. No condition making this composition the identity
has been imposed. OC45\textquotesingle s joint source/cap flow supplies the local
disc maps; its original prescribed parameter reference still supplies
the logarithm jumps to be inserted into OCP.7 and OCP.14. The lattice
calculation here removes the integral degree-two difference and gives
the displayed analytic comparison without deleting those endpoint data.

The exact program
\path{code/current_audits/audit_paper_i/iv_os45/code/k119_afterreading_exact.py}
re-inverts all twenty matrices for the four selected sites
and checks their products and primitive images.
It reconstructs the full 20-by-20 coordinate change and
inverse and the 20-by-36 mapping-cone boundary, including
the sign \(z_i=-t_i\). Together with the source, Miller,
Laurent and endpoint identities this gives 175 exact
checks, recorded in
\path{code/current_audits/audit_paper_i/iv_os45/output/k119_afterreading_exact.json}.

\par

\medskip\noindent\textbf{OCP.5. The full integral parabolic complex and its contraction}

\par\nobreak

The same actual slit-sector fillings give the based cellular complex
\[
\Lambda\xrightarrow{d^0}\bigoplus_{i=0}^4\operatorname{im}N_i
\xrightarrow{d^1}\Lambda,
\qquad d^0(v)=(N_i v)_i,\qquad
d^1((z_i)_i)=\sum_iP_i z_i.
\tag{OCP.15}
\]
The local images in this formula are the integral image lattices,
including their displayed primitive generators. The identity
\(d^1d^0=\sum_i(P_{i+1}-P_i)=0\) retains the based-cell transport.
For all four selected sites use
\[
g_0=e_w+3e_\delta,\quad g_1=e_w,\quad
g_2=e_u+3e_w+16e_\delta,\quad
g_3=e_u+8e_\delta,\quad g_4=e_u-3e_w.
\tag{OCP.16}
\]
At the selected site add \(e_\delta\) to the basis. Each unselected
image has the indicated single generator; the selected image is
exactly \(\mathbb Zg_s+\mathbb Ze_\delta\). The coefficient rows
and the primitive preimages in OCP.2-\/-3 verify these assertions over
\(\mathbb Z\). Record the five \(g_i\) coefficients first and the
selected \(e_\delta\) coefficient last, giving coordinates
\(z=(z_0,z_1,z_2,z_3,z_4,z_5)\). In these actual bases the two
matrices, independently rederived at each selected site, are
\[
D=\begin{pmatrix}
-3&-8&0&0\\
0&-1&0&0\\
-2&-3&1&0\\
-1&0&1&0\\
0&3&1&0\\
1&0&0&0
\end{pmatrix},\qquad
B=\begin{pmatrix}
0&0&0&0&0&0\\
0&0&1&-2&1&0\\
1&1&-6&9&-3&0\\
3&0&-8&8&0&1
\end{pmatrix}.
\tag{OCP.17}
\]
Thus \(D\) represents \(d^0\), and the columns of \(B\) are
\(P_i g_i\) followed by \(P_s e_\delta=e_\delta\).
The placement of the extra local generator changes with the selected
site; these matrices do not change after the stated integral basis map.

Here is a primitive preimage for every integral cocycle:
\[
Bz=0\quad\Longrightarrow\quad
z=D\begin{pmatrix}z_5\\-z_1\\z_3+z_5\\0\end{pmatrix}.
\tag{OCP.18}
\]
Indeed the three nonzero equations of \(Bz=0\) successively give
\(z_4=2z_3-z_2\),
\(z_0=3z_2-3z_3-z_1\), and
\(z_2=z_3+3z_1-z_5\). These are precisely the six rows of OCP.18.
No division occurs. For a full chain identity put
\[
Lz=(z_5,-z_1,z_3+z_5)^t,
\qquad
J=\begin{pmatrix}
-24&-8&3\\
0&0&0\\
-9&-3&1\\
0&0&0\\
10&3&-1\\
0&0&0
\end{pmatrix}.
\tag{OCP.19}
\]
Let \(D'\) be the first three columns of \(D\), and let \(B'\)
be the last three rows of \(B\), in the original \((u,w,\delta)\)
order. Direct integer multiplication gives
\[
LD'=I_3,\qquad B'J=I_3,\qquad
D'L+JB'=I_6,\qquad B'D'=0.
\tag{OCP.20}
\]
These identities are explicit splitting and contraction maps for the
actual integral complex. They prove, including all torsion information,
\[
H^0(\mathbf P^1,\mathscr H)=\mathbb Ze_\delta,
\qquad H^1(\mathbf P^1,\mathscr H)=0,
\qquad H^2(\mathbf P^1,\mathscr H)
\xrightarrow{\ \psi\ }\mathbb Z.
\tag{OCP.21}
\]
The degree-zero identification follows from the fourth zero column and
the left inverse of \(D'\). The degree-two identification follows from
the zero first row of \(B\) and its right inverse \(J\) onto the last
three coordinates. All of these are the maps for the primitive K10
fillings used in OCP.9a and OCP.15.

\par

\medskip\noindent\textbf{OCP.6. Global translating automorphisms with the retained period kernel}

\par\nobreak

For the actual Lie sheaf in OCP.11, a global section projects to a
global section of the original source-translation line
\(\mathcal O(-1)\). Its two-chart Laurent equation has only the zero
solution. The remaining physical component is a constant, so
\(H^0(\mathscr E)=\mathbb Ce_\delta\). The full off-diagonal
coefficient is still present in OCP.11-\/-13; the source component is
zero on this global-section calculation.

The actual translation exponential sequence
\(0\to\mathscr H\to\mathscr E\to\mathscr A\to0\), with the
same primitive local period kernels as OCP.15, now has global part
\[
0\longrightarrow\mathbb Ze_\delta
\longrightarrow\mathbb Ce_\delta
\longrightarrow H^0(\mathbf P^1,\mathscr A)
\longrightarrow0.
\tag{OCP.22}
\]
The last surjectivity is exactly OCP.21, with its explicit integral
contraction. In the original physical line the resulting isomorphism
is the literal map
\[
\begin{gathered}
\mathbb C/\mathbb Z\longrightarrow\mathbb C^*,
\qquad [z]\longmapsto e^{2\pi i z},\\
v_M\longmapsto e^{2\pi i z}v_M.
\end{gathered}
\tag{OCP.23}
\]
Its local inverse is \((2\pi i)^{-1}\log c\), whose two choices differ
by the retained primitive physical period. Thus every globally
extending translation is exactly such a physical scalar. The endpoint
in OCP.14 is evaluated through this same map, retaining its scalar
value and its logarithm jump; the original parameter comparison is
not replaced by a selected endpoint.

% End expanded verified cp36 input


\par\bigskip\noindent\textbf{Complete supporting proof 3: OCGC.}\par


% Expanded verified cp36 input: notes/part3_os45_cap_collisions_ocgc_proof.tex
\par

\medskip\noindent\textbf{OS45 prescribed-cap comparison: the full global maps}

\par\nobreak

The following argument proves the global source/line
comparison of OC45.6--11 and its full parameter endpoint.
It uses the local source and cap maps of OC45.1--5 and
the complete integral contraction of OCP.1--6.

\par

\medskip\noindent\textbf{OCGC.1. Original line, infinity frame, and extending translations}

\par\nobreak

Write the actual point coefficient as \(V=f_{P^8}(Z;Q)v_M\).
The translation \((Z,V)\mapsto(e^A Z,e^D V)\) acts in the original
Cartier module by
\[
v_M\longmapsto e^D\frac{f_{P^8}(Z;Q)}{f_{P^8}(e^A Z;Q)}v_M.
\tag{OCGC.1}
\]
Its inverse evaluates the reciprocal at the translated source point.
Composition cancels the intermediate actual point function. Thus its
source and target modules agree also at the point and zero-section
divisors. This is the coefficient map used in OC45.30; it is retained
before passing to a nonzero point coefficient.

The three theta reindexes at infinity give, with the full base factors,
\[
f_{Q^3b}(Z;Q)=(-Z/b)^3Q^{-6}f_b(Z;Q).
\tag{OCGC.2}
\]
Consequently the actual coefficient \(V_0=f_b(Z;Q)v_M\) changes by
\(e^{D-3A}\), while the source changes by \(e^A\). The factors
\(b^3Q^6\) remain in the forward and inverse point-coordinate maps;
they cancel in this action because the base is fixed. The period
matrix with columns \((-1,-3)^t,(3,8)^t\) has determinant one and
inverse with columns \((8,3)^t,(-3,-1)^t\). Its source valuation
row \((8,-3)\) becomes \((1,0)\). These are precisely the full
\(H=h^{-1}B^{-3},K=h^3B^8\) charts and inverse words.

At every source node the two coordinates receive reciprocal source
units; their product and the original height are fixed. At a selected
cone, a character \(Z^rV^d q^e\) receives \(e^{rA+dD}\), with its
literal reciprocal in the inverse. This respects every character
relation and every full period column. The primitive infinity chain
has the same property after the preceding determinant-one period
change. Thus the exponential used by the global comparison is the
actual extending map on the ordinary, nodal, selected, infinity, and
point/O charts, with no replacement of its units by central values.

\par

\medskip\noindent\textbf{OCGC.2. The original Lie bundle and integral comparison}

\par\nobreak

The original infinity change \(x=z^{-2}X,y=z^{-3}Y\) gives the relative
differential equality \(dx/(2y)=z\,dX/(2Y)\). Its dual frame equality
is \(e_\infty=t^{-1}e_0\). On a nodal chart the extending source
field is \(X\partial_X-Y\partial_Y\). This gives the source line
\(\mathcal O(-1)\) in the actual OS45 coordinates. The physical
vertical field is the globally defined \(\delta\). The off-diagonal
Lie coefficient is the invariant source derivative of the logarithm
of the full original point-module transition in OCGC.1-\/-2. It remains
the full function \(a(t,s)\) in
\[
\mathcal L(t,s)=\begin{pmatrix}t^{-1}&0\\a(t,s)&1\end{pmatrix}.
\tag{OCGC.3}
\]

OCP.1-\/-6 independently invert the actual twenty K10 matrices for all
four selected-zero choices. Their based matrices are
\(M_i=P_i(H_i-I)=P_{i+1}-P_i\); the explicit integral map
\(t_j=r_j-\sum_{i>j}(H_i-I)r_i\) carries the primitive raw preimages
to these transported columns. The full exterior cone sends the local
cochain \(z_i=-t_i\) to the positive retained cap-vector difference.
The parabolic complex has the displayed integral contraction
\(D'L+JB'=I_6\), \(LD'=B'J=I_3\). Thus its degree-zero kernel is
\(\mathbb Z\delta\), its degree-one group is zero integrally, and its
degree-two quotient is the original \(\mathbb Z\) measured by \(\psi\).
Both original order-one caps have \(\psi=1\). Their actual difference
therefore has the precise integral preimages in OCP.4-\/-9, including
all source-circle, physical-circle, and suspension components.

After inserting those complete period decks, the remaining Lie
cocycle is solved by OC45.53-\/-54. Substitution in OCGC.3 gives
\(t^{-1}u_\infty-u_0=f\) and
\(a u_\infty+v_\infty-v_0=g\). In particular the full term
\(a u_\infty\) is subtracted before the scalar Laurent split.
The many-disc construction first forms the actual partition cochain,
then subtracts its global Cauchy primitive. The fixed kernels give
holomorphic parameter dependence on each cut. The directed comparison
is exactly
\[
F_i=\operatorname{Exp}(-s_i),\qquad
F_i g^X_{ij}=g^Y_{ij}F_j,\qquad F_i^{-1}=\operatorname{Exp}(s_i).
\tag{OCGC.4}
\]
These are whole source and line maps through all of the preceding
charts. The base logarithms \(\log t\) and \(\log(t-d)\) retain their
individual prescribed values throughout this construction.

\par

\medskip\noindent\textbf{OCGC.3. Full parameter seam and physical scalar comparison}

\par\nobreak

On the original cap \(q_n=-s\), a positive parameter turn changes
\(h\) by \(\alpha\). For the eighteen unselected collisions the full
\(C(0,s)\) is an extending nonzero unit. For the six selected
collisions it is \(-\lambda s\beta(-s,s)\). The literal reference
\(\gamma_0+k h+r\alpha+d_\delta\delta\) therefore changes by
\(k\alpha\) and \(k\alpha+\delta\), respectively. These full vectors
are inserted into the same based integral cochain before the analytic
solve. No source or physical period has been set to zero.

A global Lie section has source component zero: the equation
\(u_0=t^{-1}u_\infty\) has no common nonzero regular Laurent power.
Its remaining component is a constant multiple of \(\delta\).
The integral degree-one vanishing just proved then identifies all
global extending translations with the exact physical action
\[
\mathbb C\delta/\mathbb Z\delta\longrightarrow\mathbb C^*,
\qquad [z\delta]\longmapsto e^{2\pi i z}.
\tag{OCGC.5}
\]
In particular ratios of the constructed comparisons are this action
on every original cap and selected fan.

For the local parameter comparisons \(F_\alpha=c_{\alpha\beta}F_\beta\),
the connection satisfies
\(\mathfrak a_\alpha-\mathfrak a_\beta=d\log c_{\alpha\beta}\).
The scalar coefficient transforms as
\(z_\alpha=c_{\alpha\beta}^{-1}z_\beta\), so its actual horizontal
equation is \(dz_\alpha=-\mathfrak a_\alpha z_\alpha\). This checks
the sign in OC45.61. The retained angular holonomy is flattened by
its explicit exponential, and the radial equation respects these
same transitions. Its antiholomorphic derivative is therefore a
single scalar form; the stated Cauchy kernel gives the global
primitive on a smaller annulus. Multiplication by its negative
exponential gives the closed holomorphic family and its reciprocal.

On the cut, the complete endpoint equality is
\[
F_1E_XF_0^{-1}=CE_Y,
\qquad C=F_1E_XF_0^{-1}E_Y^{-1}.
\tag{OCGC.6}
\]
All source/period/Cartier comparisons above prove that this particular
ratio is the physical scalar of OCGC.5. Taking the actual inverse
gives \(F_0E_X^{-1}F_1^{-1}=E_Y^{-1}C^{-1}\), without commuting any
of the other maps. Inserting the full parallel transport \(P_Y\)
then gives precisely OC45.64. The literal physical isotopy
\(e^{\theta\log C}\) extends through the same cap and cone charts.
This proves the return comparison with its supplied endpoint.

\par

\medskip\noindent\textbf{OCGC.4. Protected cap, proper disc, and the stated return}

\par\nobreak

The point \(d\) and its cap disc are chosen inside the actual nodal
coordinate neighbourhood and outside all moving critical values.
There the full \(Q,P,C\) are nonzero, and the fixed joint logarithms
give the exact plus-log map \(\mathcal F_{\log(t-d)/(2\pi i)}\)
with inverse \(\mathcal F_{-\log(t-d)/(2\pi i)}\). A base turn is the
whole original suspension \(B\), so its norm is the same positive
one and its source/line components remain those of the original flow.

At the old cap location the reference family uses the actual moving
node \(XY=t-s\), or the prescribed selected fan with this same
height. These are submersions over \(s\). The original infinity
point is \(b_0=16/27\) at the parameter centre, so its point/O module
and both full period words extend there. The protected ordinary cap,
all other simple source nodes, and all selected cone charts have the
displayed reciprocal overlap maps. Uniform contraction on the finitely
many compact pieces and the finite period representatives give the
stated compact total preimage of a closed parameter disc. Thus this
is the actual smooth proper disc family used in OC45.10.

The integral and analytic comparison carries the original positive
norm cap to that family. Its reference sphere and each prescribed
fan are carried by the retained actual ruling/Gysin/van Kampen maps.
The disc family has its identity marked endpoint; a partition of
unity applied to its local submersion lifts gives an oriented smooth
transport with inverse. OCGC.6 identifies the original return with
that transport followed by the explicitly isotopic physical scalar.
Consequently OC45.66 holds for all twenty-four stated meridians,
their reciprocals, the original order-one positive-norm reference,
and every prescribed selected fan. The radial scalar comparison
retains its argument and introduces no scalar turn. It transports
this same conclusion to every admitted original fixed scalar.

This comparison proves the scope of OC45.6-\/-11. The primitive
integral calculation, original point-module maps, complete parameter
jumps and physical endpoint morphism give the stated zero return.

% End expanded verified cp36 input


\par\bigskip\noindent\textbf{Complete supporting proof 4: ORC.}\par


% Expanded verified cp36 input: notes/part3_os45_cap_collisions_orc_proof.tex
\par

\medskip\noindent\textbf{OS45 ordinary-cap collisions: source, comparison and return}

\par\nobreak

The following argument assembles the original source,
order-one positive-norm cap, full prescribed reference,
fixed scalar and selected integral fan.
The original five K10 matrices determine its integral
comparison; the displayed source, Cartier and endpoint
maps establish its extension and return.

\par

\medskip\noindent\textbf{ORC.1. Original parameter and all twenty-four marked points}

\par\nobreak

Use the original source and original selected-node map
\[
\begin{gathered}
R=t^2+a,\quad f=R^2+t,\quad
g=R^3+\tfrac32tR+\tfrac38,\qquad
y^2=x^3-3fx+2g,\quad G=(R,\sqrt3/2),\\
a=-u/3-u^4/9,\qquad t_c=1/(2u)-u^2/3,\qquad
v^2=4u^3+3 .
\end{gathered}
\tag{ORC.1}
\]
The inverse to the marked-node substitution is the original
\(u=\xi-(t_c^2+a)\), with \(\xi\) the selected double cubic root;
the ordered normalization has \(c=v/(2u)\) and
\(Z=(S-c)/(S+c)\), \(S=c(1+Z)/(1-Z)\).
Thus the \(u,v\) labels identify the actual source branches.

The ordinary cap at \(t=0\) meets a source node exactly at
\[
16a^3+3=0,\qquad
2187-16u^3(u^3+3)^3
=-(2u^3-3)(8u^9+84u^6+342u^3+729)=0.
\tag{ORC.2}
\]
The degree-twelve polynomial is square-free, is nonzero at \(u=0\),
and has no common root with \(4u^3+3\), \(u^3+3\), or
\(16u^6+120u^3+243\). Each of its twelve roots consequently has
two distinct ordered \(v\)-values. The factor \(2u^3-3\) makes
\(t_c=0\), and gives six selected-cap points; the other factor
gives the eighteen unselected-cap points. The reciprocal branch
coordinate \(Z^{-1}\) is retained for the other \(v\)-value.

At every such point write \(a_*\) for the actual source parameter.
The finite discriminant factor and its derivatives give
\[
D=48a(t^2+a)^2+8t^3+72at+9,\quad
D_t(0,a_*)=72a_*,\quad D_a(0,a_*)=144a_*^2 .
\tag{ORC.3}
\]
The moving critical value \(s=t_n(a)\) therefore has
\(ds/da=-2a_*\). Moreover \(da/du=-(4u^3+3)/9=-v^2/9\)
is nonzero on this locus. The actual maps
\((u,v)\mapsto a\mapsto s\) are holomorphic local coordinates
with their uniquely fixed implicit inverses. A positive meridian
in the ordered source curve makes one positive \(s\)-turn.
The selected value is exactly \(s\) at the six selected-cap points.

\par

\medskip\noindent\textbf{ORC.2. Original node and complete point-line coefficient}

\par\nobreak

At \(s=0\) the source node is \((x,y,t)=(-a_*,0,0)\) and
\[
y^2=(x+a_*)^2(x-2a_*).
\tag{ORC.4}
\]
Set \(q=t-s\), choose the original branch
\(\xi(q,s)=\sqrt{f(s+q,a(s))}\) with \(\xi(0,0)=-a_*\),
and use the complete analytic unit
\[
h(q,s)=\frac{2(g(s+q,a(s))-\xi(q,s)^3)}q,\qquad h(0,0)=6a_*.
\tag{ORC.5}
\]
The numerator vanishes along \(q=0\), since it is the moving
double-root equation. Direct multiplication gives the actual maps
\[
X=y-(x-\xi)\sqrt{x+2\xi},\qquad
Y=\frac{y+(x-\xi)\sqrt{x+2\xi}}h,\qquad XY=q,
\tag{ORC.6}
\]
\[
q=XY,\quad t=s+XY,\quad y=(X+hY)/2,\qquad
2(x-\xi)\sqrt{x+2\xi}=hY-X.
\tag{ORC.7}
\]
The last equation has derivative \(2\sqrt{3\xi}\) at the node.
These maps and their inverse extend on the joint source/parameter
neighbourhood. In particular \((X,Y,s)\) are smooth total-space
coordinates with submersive projection to \(s\).

The central normalization and section are
\[
x=S^2+2a_*,\quad y=S(S^2+3a_*),\quad
c^2=-3a_*,\quad S_G=\frac{\sqrt3}{4a_*},\quad S_G^2=-a_*.
\tag{ORC.8}
\]
Thus \(c/S_G=\pm\sqrt3\) and
\(P_*=Z(G)=-2\pm\sqrt3\). Their product is one and their
absolute values differ from one. Every point \(P_*^j\) used in
the three original Miller doublings is consequently defined
and different from the zero-section coordinate one.

The original tangent normalization gives, for \(m=2,4,8\),
\[
f_m(Z)=(-2c)^{m-1}\frac{1-P^m}{(1-P)^m}
 \frac{(Z-P)^m}{(Z-P^m)(Z-1)^{m-1}}.
\tag{ORC.9}
\]
Its divisor is \(m[G]-[mG]-(m-1)[O]\); substituting
\(S=c(1+Z)/(1-Z)\) shows the principal coefficient at \(O\)
is exactly \(S^{m-1}\). These two checks agree with the
successive original tangent quotients. For the actual direction
\(T=-G\), the full chronological coefficient is
\[
B_{8,1}=f_2^8\frac{f_8(Z/P)}{f_8(Z)}
 =(-2S_G)^8\frac{(Z-P)(Z-P^8)}{(Z-P^9)(Z-1)}.
\tag{ORC.10}
\]
With the actual point function \(f_b=(1-Z/b)/(1-Z)\),
\[
B_{8,1}(Z)\frac{f_{P^8}(Z/P)}{f_{P^8}(Z)}
=(-2S_G)^8=\frac{81}{256a_*^8}=\frac9{a_*^2}.
\tag{ORC.11}
\]
The equality cancels the four displayed source and target Cartier
factors; their point/O modules remain the modules of the map.
The inverse evaluates the reciprocal at the inverse source point.

Retain the complete inverse-Tate functions \(Q=q\kappa(q,s)\),
\(P=P(q,s)\), and the actual Miller/theta coefficient \(\beta(q,s)\),
fixed at the original regular reference. Their central values are
\[
\kappa(0,0)=-\frac{h(0,0)}{1728\xi(0,0)^3}
 =\frac1{288a_*^2},\qquad
\beta(0,0)=C_*\frac9{a_*^2}.
\tag{ORC.12}
\]
Here \(C_*\) is the original fixed Hom normalization. The full
period arrows, before any local logarithm, are
\[
h_{\mathrm{per}}:(Z,V)\longmapsto(QZ,P^8V),\qquad
B:(Z,V)\longmapsto(P^{-1}Z,CV),\quad
C=\lambda(t-t_c)\beta .
\tag{ORC.13}
\]
All higher coefficients of these functions occur in the maps below.

\par

\medskip\noindent\textbf{ORC.3. Extending translations and the infinity component map}

\par\nobreak

A Lie translation with complete analytic coefficients \(A,D\) is
\[
\mathcal T_{A,D}(Z,V,t,s)=(e^AZ,e^DV,t,s),
\quad
\mathcal T_{A,D}^{-1}=\mathcal T_{-A,-D}.
\tag{ORC.14}
\]
The base is fixed. On the original point module its coefficient is
\[
e^D\frac{f_{P^8}(Z;Q)}{f_{P^8}(e^AZ;Q)}.
\tag{ORC.15}
\]
Composition cancels the middle point function at the actual
intermediate source point. The reciprocal there proves the
inverse identity. At a zero or pole the numerator and denominator
are the source and target Cartier generators, so the same formula
is the module isomorphism through their graphs.

On a source chain, \(X_i=Z/Q^i,\ Y_i=Q^{i+1}/Z\) receive the
units \(e^A,e^{-A}\). Their product \(Q\), and hence the original
height \(q\) under \(Q=q\kappa\), is fixed. At a selected fan cone,
the actual character \(Z^rV^d q^e\) receives \(e^{rA+dD}\).
This proves the map, its inverse, and every monomial overlap on
each prescribed unimodular fan.

At infinity keep \(z=1/t\), the original minimal source coordinates,
and the full functions
\[
Q_\infty=z^8\kappa_\infty,\qquad
P_\infty=z^3\pi,\qquad
b_0=P_\infty^8/Q_\infty^3,\qquad
b_0(0,s)=-1/(9a(s)^3).
\tag{ORC.16}
\]
The complete theta identity is
\[
f_{Q_\infty^3b_0}(Z;Q_\infty)
 =(-Z/b_0)^3Q_\infty^{-6}f_{b_0}(Z;Q_\infty).
\tag{ORC.17}
\]
Writing \(V_0=f_{b_0}v_M\) therefore gives
\((Z,V_0)\mapsto(e^AZ,e^{D-3A}V_0)\).
Both base factors \(b_0^3Q_\infty^6\) occur in the reciprocal
coordinate map; fixing the base cancels them in this action.

The original source valuations of \(h_{\mathrm{per}},B\) are
\((8,-3)\). Their actual period words and inverse are
\[
H=h_{\mathrm{per}}^{-1}B^{-3},\quad
K=h_{\mathrm{per}}^3B^8,\qquad
h_{\mathrm{per}}=H^8K^3,\quad B=H^{-3}K^{-1}.
\tag{ORC.18}
\]
The integer matrix
\(\left(\begin{smallmatrix}-1&3\\-3&8\end{smallmatrix}\right)\)
has inverse
\(\left(\begin{smallmatrix}8&-3\\3&-1\end{smallmatrix}\right)\);
the row \((8,-3)\) becomes \((1,0)\).
The line multiplier of \(K\) is the full unit \(\lambda^8\)
times its original analytic coefficient. In particular every
component shift has its explicit inverse period word: a source
order \(n\) is removed by \(H^{-n}\). These words are evaluated
on the original source and line, including their unit factors.
The rank-one chain action of (ORC.14) consequently extends
through every original infinity component and descends to the
same compact quotient. This proves the component transfer used
by the exponential, with its inverse.

\par

\medskip\noindent\textbf{ORC.4. The actual local exponential and original Lie bundle}

\par\nobreak

Let \(\mathscr A\) be the sheaf of extending fibre translations
in the just displayed charts, \(\mathscr E\) their Lie bundle,
and \(\mathscr H\) the period kernel. The local inverse to
\(\operatorname{Exp}\) is obtained as follows. On an unselected
node, remove the integral source order by the period
\(h_{\mathrm{per}}\), whose source valuation is one. At the
selected filling the two period valuations are exactly
\((1,0),(0,1)\); their integral inverse removes the source
and line orders. At infinity use (ORC.18) to remove the source
order. The remaining line multiplier has order zero because
the map and its inverse preserve the nonzero line quotient.
In each chart the remaining multipliers are nonvanishing
analytic units; their analytic logarithms give \(A,D\).
The differences of two logarithms are precisely the retained
angular periods and the extending zero-radial-order period
words. All removals are period maps on the quotient, and
their inverse words were specified above.

The radial coefficients of a local full period vanish exactly
when the original monodromy fixes that period. Thus
\(\mathscr H=j_*\Lambda\), with stalk \(\ker(H_i-I)\) at the
five source fillings, and the explicit maps give
\[
0\longrightarrow\mathscr H\longrightarrow\mathscr E
 \xrightarrow{\operatorname{Exp}}\mathscr A\longrightarrow0.
\tag{ORC.19}
\]
On a slit source disc the annular complex is
\(\Lambda\xrightarrow{H_i-I}\Lambda\). Attaching its centre
keeps its invariant zero-cochains; the same annular
coinvariant term occurs with the opposite boundary in the
disc/annulus mapping cone and cancels by its identity map.
This gives the invariant stalk and zero positive local
cohomology used in the global based complex below.

The original coordinate equality
\[
\frac{dx}{2y}=z\,\frac{dX}{2Y},\qquad
e_\infty=t^{-1}e_0
\tag{ORC.20}
\]
identifies the source Lie line with \(\mathcal O(-1)\).
At a source node its vector is
\(X_i\partial_{X_i}-Y_i\partial_{Y_i}\).
The physical vector \(\delta\) has literal flow
\(v_M\mapsto e^{2\pi iw}v_M\). Differentiating the entire
point-module transitions (ORC.15),(ORC.17) gives
\[
0\longrightarrow\mathcal O\delta\longrightarrow\mathscr E
 \longrightarrow\mathcal O(-1)\longrightarrow0,\qquad
\mathcal L(t,s)=
\begin{pmatrix}t^{-1}&0\\ a(t,s)&1\end{pmatrix}.
\tag{ORC.21}
\]
The coefficient \(a(t,s)\) is this full logarithmic derivative.
No truncation of it occurs in the cochain solution.

\par

\medskip\noindent\textbf{ORC.5. Literal integral contraction of the five source fillings}

\par\nobreak

Take the five original K10 matrices in the order
\((\psi,u,w,\delta)\) and set
\[
H_i=T_i^{-t},\quad N_i=H_i-I,\quad
P_0=I,\quad P_i=H_0\cdots H_{i-1},\quad
M_i=P_iN_i=P_{i+1}-P_i.
\tag{ORC.22}
\]
The original product is \(H_0\cdots H_4=I\).
For the selected site \(\sigma=1,2,3,4\), respectively, use
\[
d_1=e_\psi,\quad d_2=e_\psi+2e_w,\quad
d_3=e_\psi+e_w,\quad d_4=e_\psi .
\tag{ORC.23}
\]
Literal inversion of the original matrices, with their selected
central entries, gives
\[
N_\sigma d_\sigma=e_\delta,\quad
N_1(-e_u)=e_w,\quad
N_3e_w=e_u+8e_\delta.
\tag{ORC.24}
\]
For every retained vector \(v=Ae_u+Be_w+Ce_\delta\), set
\[
r_1=-Be_u,\quad r_3=Ae_w,\quad
r_\sigma\mathrel{+}=(C-8A)d_\sigma,\qquad
t_j=r_j-\sum_{i>j}N_ir_i.
\tag{ORC.25}
\]
Repeated indices receive both terms. The inverse is the
descending recursion \(r_j=t_j+\sum_{i>j}N_ir_i\).
The telescoping equality \(\sum_{j<i}M_j=P_i-I\) gives
\(\sum_jM_jt_j=\sum_iN_ir_i=v\).
Thus these are inverse integral based-cell coordinate maps.

The oriented cone boundary is
\[
b_i=p_i-N_iz_i\quad(i<4),\qquad
b_4=-H_4\sum_{i<4}P_ip_i-N_4z_4.
\tag{ORC.26}
\]
Its coordinate \(v_*=\sum_iP_ib_i\) cancels every \(p_i\)
and equals \(-\sum_iM_iz_i\); hence the positive \(v\)
uses precisely \(z_i=-t_i\). This keeps the original
suspension terms in \(d_\sigma\) with their whole period
arrows, including at both ends of the corresponding cut.

After the local disc/annulus cancellation in (ORC.19), the
remaining integral complex is
\[
\Lambda\ \xrightarrow{D}\ \mathbb Z^6\
 \xrightarrow{B}\ \Lambda .
\tag{ORC.27}
\]
The first five local generators are
\[
e_w+3e_\delta,\ e_w,\ e_u+3e_w+16e_\delta,\
e_u+8e_\delta,\ e_u-3e_w;
\]
the sixth is \(e_\delta\) at the selected site. Their primitive
preimages in (ORC.24) give the full image lattices. In those
actual bases, for each of the four selected sites,
\[
D=\begin{pmatrix}
-3&-8&0&0\\0&-1&0&0\\-2&-3&1&0\\
-1&0&1&0\\0&3&1&0\\1&0&0&0
\end{pmatrix},\qquad
B=\begin{pmatrix}
0&0&0&0&0&0\\0&0&1&-2&1&0\\
1&1&-6&9&-3&0\\3&0&-8&8&0&1
\end{pmatrix}.
\tag{ORC.28}
\]
Write \(D'\) for the first three columns and \(B'\) for the
last three rows. The exact splitting maps are
\[
Lz=(z_5,-z_1,z_3+z_5)^t,\qquad
J=\begin{pmatrix}
-24&-8&3\\0&0&0\\-9&-3&1\\0&0&0\\
10&3&-1\\0&0&0
\end{pmatrix},
\tag{ORC.29}
\]
\[
LD'=I_3,\quad B'J=I_3,\quad
D'L+JB'=I_6,\quad B'D'=0.
\tag{ORC.30}
\]
Every map is integral. These identities supply the inverse
on each image and the contraction of degree one, including
all torsion information. They prove
\[
H^0(\mathscr H)=\mathbb Z\delta,\qquad
H^1(\mathscr H)=0,\qquad
H^2(\mathscr H)\xrightarrow{\psi}\mathbb Z ,
\tag{ORC.31}
\]
where the last map has section \(1\mapsto[e_\psi]\).
Smooth cap discs add their transported full period vector
to \(v_*\). The two order-one positive-norm caps considered
below therefore differ by exactly the vector \(v\) in
(ORC.25), retaining all three of its integer coordinates.

\par

\medskip\noindent\textbf{ORC.6. Full analytic map between the two cap constructions}

\par\nobreak

Take the original cap at \(t=0\) and a protected ordinary cap
at \(t=d\). Use a specified based path on a cut parameter
disc to transport their full lattice markings. Their actual
clutch logarithms use \(\log t\) and \(\log(t-d)\).
Their difference has a period-valued two-cocycle with
\(\psi=1-1=0\). Subtract the explicit cochain
(ORC.25)-\/-(ORC.26), evaluating every inserted period using
the original source, line and scalar functions. The two
branches on a slit annulus differ by the cancelled
\(N_iz_i\), so the remaining Lie cocycle is single valued.

Here is its analytic solution with the full off-diagonal
coefficient. For an overlap vector \((f,g)\) in the frames
(ORC.21), let its Laurent coefficients be the fixed-contour
Cauchy integrals and put
\[
u_0=-\sum_{j\geq0}f_jt^j,\qquad
u_\infty=\sum_{j<0}f_jt^{j+1},\qquad
G=g-a(t,s)u_\infty ,
\tag{ORC.32}
\]
\[
v_0=-\sum_{j\geq0}G_jt^j,\qquad
v_\infty=\sum_{j<0}G_jt^j.
\tag{ORC.33}
\]
Direct substitution gives
\(\mathcal L(u_\infty,v_\infty)^t-(u_0,v_0)^t=(f,g)^t\).
Both infinity components are holomorphic in \(1/t\);
the \(j=-1\) source coefficient is an allowed constant.

On the full many-disc cover first form the partition cochain
\(s_i^{\mathrm{sm}}=\sum_k\rho_k\theta_{ik}\).
The cocycle equation gives
\(s_i^{\mathrm{sm}}-\mathcal L_{ij}s_j^{\mathrm{sm}}=\theta_{ij}\).
Its antiholomorphic derivative is a global \(\mathscr E\)-valued
form. Componentwise Cauchy primitives on the two affine
charts have holomorphic difference; (ORC.32)-\/-(ORC.33)
remove that difference. Subtracting this global primitive
gives holomorphic \(s_i\) with the same cocycle.
Fixed partitions and kernels retain holomorphic dependence
on the cut parameter. This construction also proves
\(H^1(\mathscr E)=0\) with its actual transition.

For directed overlaps \(g^X_{ij}=t_{\theta_{ij}}g^Y_{ij}\),
the resulting full maps are
\[
F_i=\operatorname{Exp}(-s_i),\qquad
F_i^{-1}=\operatorname{Exp}(s_i),\qquad
F_i g^X_{ij}=g^Y_{ij}F_j .
\tag{ORC.34}
\]
Equations (ORC.14)-\/-(ORC.18) give their literal coefficients
and reciprocals in the original point/O, source-chain,
infinity and selected-fan coordinates. On each smooth cap
they are the same exponential before its plus-log gluing.
Thus the comparison carries the complete original cap
reference and its source/physical integers.

\par

\medskip\noindent\textbf{ORC.7. Both parameter jumps and the global comparison group}

\par\nobreak

At the original cap \(q=-s\), write the actual period on a cut as
\[
\gamma_{\mathrm{ref}}=
\frac{(-\log P,\log C)}{2\pi i}
+k\,\frac{(\log Q,8\log P)}{2\pi i}
+r\alpha+d_\delta\delta .
\tag{ORC.35}
\]
All coefficients \(k,r,d_\delta\) are the original integers
in the based marking. A positive \(s\)-turn adds \(\alpha\)
to the second displayed full period. At an unselected-cap
point \(C(0,s)\) is an extending unit. At a selected-cap
point it is exactly \(-\lambda s\beta(-s,s)\). Therefore
\[
\Delta\gamma_{\mathrm{ref}}=k\alpha
\quad(18\text{ points}),\qquad
\Delta\gamma_{\mathrm{ref}}=k\alpha+\delta
\quad(6\text{ points}).
\tag{ORC.36}
\]
The logarithms of every extending unit have their full
continued values. These vectors have zero \(\psi\)-coordinate;
their based transports are inserted into (ORC.25)-\/-(ORC.34)
at the final parameter cut. This retains the physical jump
at the six selected-cap points.

A global source Lie section satisfies \(u_0=t^{-1}u_\infty\).
Regular powers at zero and infinity force it to vanish.
The remaining global Lie section is a constant physical
vector. The exact exponential (ORC.19), together with the
integral degree-one contraction (ORC.30), gives
\[
H^0(\mathscr A)=\mathbb C\delta/\mathbb Z\delta
 \xrightarrow{\ \sim\ }\mathbb C^*,\qquad
[w\delta]\longmapsto(v_M\mapsto e^{2\pi iw}v_M).
\tag{ORC.37}
\]
The local inverse is \((2\pi i)^{-1}\log c\), and two
choices differ by the retained primitive period.
Consequently every ratio of the constructed global
comparisons is this literal physical scalar on every
original chart.

\par

\medskip\noindent\textbf{ORC.8. Closed parameter comparison and directed endpoint}

\par\nobreak

For local parameter comparisons \(F_\alpha=c_{\alpha\beta}F_\beta\),
all \(c_{\alpha\beta}\) are the scalars in (ORC.37). Choose a
partition and form
\(\mathfrak a_\alpha=\sum_\beta\rho_\beta\,d\log c_{\alpha\beta}\).
Then
\(\mathfrak a_\alpha-\mathfrak a_\beta=d\log c_{\alpha\beta}\).
The scalar coefficients of a comparison transform by
\(z_\alpha=c_{\alpha\beta}^{-1}z_\beta\), so their horizontal
equation is exactly
\[
dz_\alpha=-\mathfrak a_\alpha z_\alpha,\qquad
z_\alpha(\rho,\theta)=z_\alpha(\rho_0,\theta)
\exp\!\left(-\int_{\rho_0}^{\rho}
\mathfrak a_\alpha(\partial_r)\,dr\right).
\tag{ORC.38}
\]
Angular transport has a nonzero scalar holonomy \(c\).
Multiplication by \(e^{-\chi(\theta)\log c}\), with
\(\chi(\theta+1)=\chi(\theta)+1\) and constant endpoint
collars, makes the comparison periodic with all derivatives.
The radial equation extends it to a thin annulus.
Its antiholomorphic derivative is a single scalar form
\(\bar\partial F^{\mathrm{sm}}=\eta F^{\mathrm{sm}}\).
After a compactly supported extension, the Cauchy integral
\[
u(s)=\frac1\pi\int_{\mathbb C}\frac{\eta_{\bar w}(w)}{s-w}\,dA(w)
\tag{ORC.39}
\]
satisfies \(\bar\partial u=\eta\) on a smaller annulus.
Thus \(e^{-u}F^{\mathrm{sm}}\) is a nonzero closed
holomorphic comparison, with its explicit reciprocal.

The endpoint direction is also visible without choosing
this closed section. Let \(F_\theta:X_\theta\to Y_\theta\)
be (ORC.34), and let \(E_X:X_0\to X_1\),
\(E_Y:Y_0\to Y_1\) be the full forward marked endpoints.
All continued periods in (ORC.36) are included. The exact
ratio and its inverse equation are
\[
C=F_1E_XF_0^{-1}E_Y^{-1}\in\mathbb C^*_{\mathrm{physical}},
\qquad F_0E_X^{-1}F_1^{-1}=E_Y^{-1}C^{-1}.
\tag{ORC.40}
\]
For actual parallel transport \(P_Y:Y_0\to Y_1\),
\[
F_0\bigl(E_X^{-1}F_1^{-1}P_YF_0\bigr)F_0^{-1}
=E_Y^{-1}C^{-1}P_Y.
\tag{ORC.41}
\]
This follows by composition in the written order. No
commutation with a source endpoint or transport is used.
The scalar factor has the physical isotopy
\(e^{\theta\log C}\), whose maps and inverses extend on
every original cap, line, infinity and selected fan.

\par

\medskip\noindent\textbf{ORC.9. Actual protected-cap family across the parameter centre}

\par\nobreak

Choose a fixed \(d\ne0\) in the original nodal coordinate
neighbourhood and a cap disc of radius less than \(|d|/4\).
A sufficiently small closed \(s\)-disc keeps it disjoint
from all moving critical values and the selected value.
There the complete \(Q,P,C\) are nonzero, and their
analytic logarithms are fixed at the same reference.
The actual flow and cap maps are
\[
\mathcal F_w(Z,V)=
(e^{-w\log P}Z,e^{w\log C}V),\quad
\mathcal F_w^{-1}=\mathcal F_{-w},\qquad
\Psi_d=\mathcal F_{\log(t-d)/(2\pi i)}.
\tag{ORC.42}
\]
A branch change gives the whole \(B\)-period. The maps
therefore descend to the compact quotient, and their
reciprocal is \(\mathcal F_{-\log(t-d)/(2\pi i)}\).
This is the same positive norm one, with all other
period components recorded by the full logarithms.

At the old cap location the family \(Y\) uses the actual
source node (ORC.6)-\/-(ORC.7), or the prescribed selected
fan with \(q=t-s\). The source coordinates \((X,Y,s)\)
and the unimodular cone coordinates with \(s\) give
submersions. At infinity (ORC.16) has
\(b_0(0,0)=16/27\), so its point/O modules and full
period words extend. The other source critical values
remain distinct and simple by the original discriminant
and the disjoint polynomial loci in ORC.1. Their original
charts and full unit overlaps extend on the same disc.

Choose a common strictly contracting scalar for these
compact source pieces. The contracting line fundamental
annuli are compact. At infinity the actual \(H,K\) words
give a compact fundamental set for the primitive chain
quotient; the original eight components are carried by
their inverse words. At the selected filling the two
primitive valuation columns and finitely many cones
modulo their periods give a compact fundamental set.
The protected smooth cap and the remaining finite
source pieces are compact over their closed base discs.
All unit multipliers are bounded away from zero on
the closed parameter disc. Gluing these finitely many
compact pieces by the stated collar maps and their
inverses gives a Hausdorff compact preimage of that disc.
The local submersions agree on the same maps. This
establishes a smooth proper family \(Y\) across \(s=0\).

On a reference parameter, (ORC.34) identifies \(Y\)
with the original K10 sphere with its full cap reference.
Each prescribed fan is carried from that sphere by the
retained finite canonical fan word, the actual common
blowup and two-ruling maps, and their primitive Gysin
and van Kampen maps in F35/O45P.6. Thus the disc
family has the required oriented sphere fibre.
Partitions of unity on the displayed submersion charts
give a global horizontal lift; compactness gives its
flow and inverse for every radial and angular path.
A radial contraction of the parameter circle supplies
the isotopy between its return and the identity.

For an original admitted scalar outside the common
small disc, use the actual path
\(\lambda_\tau=((1-\tau)+\tau\varepsilon)\lambda_0\)
with \(0<\varepsilon<1\). The complete suspension
coefficient is linear in \(\lambda\), so this path
remains contracting and has zero argument change.
The same source, point-line, cap and fan arrows give
the compact cylinder and its inverse transports.
This conjugates the marked returns to the common-disc
calculation while preserving the original endpoints.

\par

\medskip\noindent\textbf{ORC.10. Original-reference return and precise scope}

\par\nobreak

The full maps of \(Y\) are defined on the parameter
disc, so its forward endpoint around the circle is
the identity. Its actual return is isotopic to the
identity by the flow just constructed.
Equation (ORC.41) transfers this return to the original
family \(X\), and (ORC.40) removes the remaining
physical scalar by its explicit isotopy. Therefore
\[
[\Phi]=0\quad\text{in }\Theta_7=\mathbb Z/28
\tag{ORC.43}
\]
for each of the twenty-four ordinary-cap collision
meridians of the original ordered OS45 source cover,
with its original order-one positive-norm cap, complete
prescribed reference, every prescribed lattice-basis
selected fan, and original admitted fixed scalar.
Each inverse meridian uses the reciprocal source,
line, cap, endpoint and transport maps and has the
same zero class.

For the eighteen unselected-cap points the original
particular period also gives the local enlarged-cap
construction: the middle ordinary region maps by
\(\Psi_0^{-1}\), the small cap and exterior map by
their identities, and both collar equations are
\(\Psi_0^{-1}\Psi_0=1\). Its inverse is obtained
piecewise with the same full flow. The original
reference assertion above additionally carries every
retained integer and the six selected-cap jumps.
Part3.4 remains active on the separately labelled
original-source generators outside this completed
ordinary-cap domain.

% End expanded verified cp36 input



\par\medskip\noindent\textbf{Source maps and exact calculation.}\par
The source, point line, five marked monodromy matrices and primitive
fillings are those of \cref{night:P3K.116}; the prescribed-fan sphere
transport uses the common blowup and ruling maps of \cref{night:P3F.35}.
The four complete arguments above retain the source, Cartier, integral,
analytic, proper-family and directed endpoint maps.

The program
\path{audit_paper_i/iv_os45/code/k119_afterreading_exact.py}, also
included under \path{code/current_audits/}, reconstructs the four
selected-zero choices and all twenty inverse-transpose matrices.
It computes the full $20\times20$ raw/based maps and their integral
inverses, the $20\times36$ cone boundary, the primitive image
coordinates and the degree-one contraction. Its 175 exact checks
also verify the original source/node and chronological Miller
identities, the full Laurent solve and the noncommutative endpoint
word. The output is
\path{audit_paper_i/iv_os45/output/k119_afterreading_exact.json}.

% End expanded verified cp36 input

\end{proof}


% END CONTINUATION_NIGHT36_RECEIVING



\subsection{The original M2 selected-I3 return in the actual fan \texorpdfstring{\(T_{-1}\)}{T-1}}
